\subsection{Spectrum and resolvent}

Lemma 4. --- Let E be a complex Banach space and let u be a closed injective partial operator on E such that \(u^{-1} \in \mathcal{L}(\mathrm{E})\) . Let \(v \in \mathcal{L}(\mathrm{E})\)

such that \(\|v\| < \|u^{-1}\|^{-1}\) . Then the partial operator \(u + v\) is injective, we have \((u + v)^{-1} \in \mathcal{L}(\mathrm{E})\) and

\[
(u + v) ^ {- 1} = u ^ {- 1} \circ \sum_ {k = 0} ^ {+ \infty} (- v u ^ {- 1}) ^ {k},\tag{6}
\]

where the series is absolutely convergent in \(\mathcal{L}(\mathrm{E})\) . Moreover, we have

\[
\| (u + v) ^ {- 1} \| \leqslant \frac {\| u ^ {- 1} \|}{1 - \| v \| \| u ^ {- 1} \|}.
\]

Since \(\|v\|\|u^{-1}\|<1\) , the series with general term \((-vu^{-1})^{k}\) is absolutely convergent in \(\mathcal{L}(\mathrm{E})\) and its sum is the inverse of the endomorphism \(1_{E}+vu^{-1}\) (prop. 2 of I, p. 22). Consequently, the partial operator \((1_{\mathrm{E}}+vu^{-1})\circ u=u+v\) (IV, p. 226, rem. vi) is injective and the reciprocal partial operator \((u+v)^{-1}=u^{-1}\circ(1_{\mathrm{E}}+vu^{-1})^{-1}\) (IV, p. 226, rem. iii) belongs to \(\mathcal{L}(\mathrm{E})\) . Since

\[
\left\| \left(1 _ {\mathrm{E}} + v u ^ {- 1}\right) ^ {- 1} \right\| \leqslant \sum_ {k = 0} ^ {+ \infty} (\left\| v \right\| \left\| u ^ {- 1} \right\|) ^ {k} = \frac {1}{1 - \left\| v \right\| \left\| u ^ {- 1} \right\|},
\]

the lemma is proved.

DEFINITION 8. --- Let u be a closed operator with dense domain on a complex Banach space E. The resolvent set of u is the set of complex numbers λ such that the partial operator \(\lambda 1_{\mathrm{E}} - u\) is injective and its inverse \((\lambda 1_{\mathrm{E}} - u)^{-1}\) belongs to \(\mathcal{L}(\mathrm{E})\) .

The spectrum of u, denoted \(\mathrm{Sp}(u)\) , is the complement of the resolvent set in C.

If \(\lambda \in \mathbf{C} - \mathrm{Sp}(u)\) , we denote by \(\mathrm{R}(u,\lambda) \in \mathcal{L}(\mathrm{E})\) the inverse of \(\lambda 1_{\mathrm{E}} - u\) . The map from \(\mathbf{C} - \mathrm{Sp}(u)\) into \(\mathcal{L}(\mathrm{E})\) that sends \(\lambda\) to \(\mathrm{R}(u,\lambda)\) is called the resolvent of \(u\) .

Remarks. --- Let E be a complex Banach space and let u be a closed operator with dense domain on E.

\begin{enumerate}
\def\labelenumi{\arabic{enumi})}
\item
  If \(u \in \mathcal{L}(\mathrm{E})\) , its spectrum coincides with the spectrum of the element \(u\) of the algebra \(\mathcal{L}(\mathrm{E})\) (I, p. 2, def. 1).
\item
  Let D be the domain of \(u\) . For every \(\lambda \in \mathbf{C} - \mathrm{Sp}(u)\) , we have
\end{enumerate}

\[
1 _ {\mathrm{E}} = (\lambda 1 _ {\mathrm{E}} - u) \circ \mathrm{R} (u, \lambda), \quad 1 _ {\mathrm{D}} = \mathrm{R} (u, \lambda) \circ (\lambda 1 _ {\mathrm{E}} - u).\tag{7}
\]

Moreover, for \(\lambda_{1}\) and \(\lambda_{2}\) in the resolvent set of u, we have

\[
\mathrm{R} (u, \lambda_ {1}) - \mathrm{R} (u, \lambda_ {2}) = (\lambda_ {2} - \lambda_ {1}) \mathrm{R} (u, \lambda_ {2}) \circ \mathrm{R} (u, \lambda_ {1}),\tag{8}
\]

\[
\mathrm{R} (u, \lambda_ {1}) \circ \mathrm{R} (u, \lambda_ {2}) = \mathrm{R} (u, \lambda_ {2}) \circ \mathrm{R} (u, \lambda_ {1}),\tag{9}
\]

Indeed, one has

\[
\mathrm{R} (u, \lambda_ {1}) - \mathrm{R} (u, \lambda_ {2}) = \mathrm{R} (u, \lambda_ {1}) \circ 1 _ {\mathrm{E}} - 1 _ {\mathrm{D}} \circ \mathrm{R} (u, \lambda_ {2}).
\]

Since \(1_{\mathrm{E}} = (\lambda_21_{\mathrm{E}} - u)\circ \mathrm{R}(u,\lambda_2)\) and \(1_{\mathrm{D}} = \mathrm{R}(u,\lambda_1)\circ (\lambda_11_{\mathrm{E}} - u)\) , we obtain

\[
\begin{array}{c} \mathrm{R} (u, \lambda_ {1}) - \mathrm{R} (u, \lambda_ {2}) = \lambda_ {2} \mathrm{R} (u, \lambda_ {1}) \circ \mathrm{R} (u, \lambda_ {2}) \\ - \mathrm{R} (u, \lambda_ {1}) \circ u \circ \mathrm{R} (u, \lambda_ {2}) - \lambda_ {1} \mathrm{R} (u, \lambda_ {1}) \circ \mathrm{R} (u, \lambda_ {2}) \\ + \mathrm{R} (u, \lambda_ {1}) \circ u \circ \mathrm{R} (u, \lambda_ {2}), \end{array}
\]

whence the first formula. Exchanging the roles of \(\lambda_{1}\) and \(\lambda_{2}\) we deduce the second formula.

\begin{enumerate}
\def\labelenumi{\arabic{enumi})}
\setcounter{enumi}{2}
\item
  Let \(\lambda \in \mathbf{C}\) . By the closed graph theorem (EVT, I, p. 19, cor. 5), if the linear map from \(\mathrm{dom}(u)\) to E defined by \(x \mapsto (\lambda 1_{\mathrm{E}} - u)(x)\) is bijective, then its inverse, whose graph is closed, is continuous from E to E. Therefore \(\lambda\) belongs to the resolvent set of \(u\) if and only if the map \(x \mapsto (\lambda 1_{\mathrm{E}} - u)(x)\) from \(\mathrm{dom}(u)\) into E is bijective.
\item
  If \(\lambda\) belongs to the spectrum of u, one of the following properties holds:
\end{enumerate}

\begin{enumerate}
\def\labelenumi{(\roman{enumi})}
\item
  The kernel of \(\lambda 1_{\mathrm{E}} - u\) is not reduced to 0; one then says that \(\lambda\) is an eigenvalue of \(u\) , and that the dimension of \(\operatorname{Ker}(\lambda 1_{\mathrm{E}} - u)\) is its multiplicity;
\item
  The partial operator \(\lambda 1_{\mathrm{E}} - u\) is injective and its image is not dense in E; one says that \(\lambda\) belongs to the residual spectrum of \(u\) ;
\item
  The partial operator \(\lambda 1_{\mathrm{E}} - u\) is injective, its image is dense in E, but \(\lambda 1_{\mathrm{E}} - u\) is not surjective; one says that \(\lambda\) belongs to the continuous spectrum of \(u\) .
\end{enumerate}

\begin{enumerate}
\def\labelenumi{\arabic{enumi})}
\setcounter{enumi}{4}
\tightlist
\item
  Let \(\lambda\) be a complex number belonging to the resolvent set of \(u\) . The resolvent \(\mathrm{R}(u,\lambda)\) is an injective linear map from E into E; its image is the domain of \(u\) and \(u = \lambda 1_{\mathrm{E}} - \mathrm{R}(u,\lambda)^{-1}\) (cf. IV, p. 226). Conversely, this property characterizes the resolvent set and the resolvent. Precisely, let \(\lambda \in \mathbf{C}\) ; if there exists a continuous injective linear map w from E into E such that \(u = \lambda 1_{E} - w^{-1}\) , then \(\lambda\) belongs to the resolvent set of u and \(w = \mathrm{R}(u, \lambda)\) .
\end{enumerate}

In particular, if v is a closed densely defined operator on E, and if \(\lambda \in C\) is a complex number not belonging to \(\operatorname{Sp}(u) \cup \operatorname{Sp}(v)\) , then the equality \(\operatorname{R}(u, \lambda) = \operatorname{R}(v, \lambda)\) implies that u = v.

\begin{enumerate}
\def\labelenumi{\arabic{enumi})}
\setcounter{enumi}{5}
\tightlist
\item
  We define the spectrum of a closable partial operator \(u\) as the spectrum of its closure.
\end{enumerate}

There exist closed operators whose spectrum is empty, or whose spectrum is equal to C (Exercise 12 of IV, p. 347).

Let E be a complex Banach space and u a closed densely defined operator on E. If F is a complex Banach space and if \(v: E \to F\) is an isomorphism, then we have \(\mathrm{Sp}(v \circ u \circ v^{-1}) = \mathrm{Sp}(u)\) and \(\mathrm{R}(v \circ u \circ v^{-1}, \lambda) = v \circ \mathrm{R}(u, \lambda) \circ v^{-1}\) for all \(\lambda \notin \mathrm{Sp}(u)\) .

PROPOSITION 14. --- Let E be a complex Banach space. Let u be a closed operator with dense domain on E and let U be its resolvent set. a) For every \(\lambda\in\mathrm{U}\) the open disk in C with center \(\lambda\) and radius \(\|\mathrm{R}(u,\lambda)\|^{-1}\) is contained in U;

\begin{enumerate}
\def\labelenumi{\alph{enumi})}
\setcounter{enumi}{1}
\item
  The set U is open in C;
\item
  Suppose that \(\operatorname{Sp}(u)\) is not empty. Let \(\lambda \in \mathrm{U}\) and let \(\delta\) be the distance in \(\mathbf{C}\) of \(\lambda\) to the spectrum of \(u\) . We have \(\delta > 0\) and \(\| \mathrm{R}(u, \lambda) \| \geqslant 1 / \delta\) ;
\item
  The map \(\lambda \mapsto \mathrm{R}(u,\lambda)\) is a holomorphic map from U to \(\mathcal{L}(\mathrm{E})\) . For every integer \(k\in \mathbf{N}\) and every \(\lambda \in \mathrm{U}\) , we have
\end{enumerate}

\[
\frac {\partial^ {k}}{\partial \lambda^ {k}} \mathrm{R} (u, \lambda) = (- 1) ^ {k} k! \mathrm{R} (u, \lambda) ^ {k + 1}.
\]

Let \(\lambda \in \mathrm{U}\) . For every \(\mu \in \mathbf{C}\) such that \(\| (\mu -\lambda)1_{\mathrm{E}}\| < \| \mathrm{R}(u,\lambda)\|^{-1}\) , Lemma 4 applied to the injective partial operator \(\lambda 1_{\mathrm{E}} - u\) and with \(v = (\mu -\lambda)1_{\mathrm{E}}\) implies that the partial operator \(\mu 1_{\mathrm{E}} - u = \lambda 1_{\mathrm{E}} - u + v\) is injective and has a continuous inverse. This implies \(a\) ). By loc. cit., we also have the formula

\[
\mathrm{R} (u, \mu) = \mathrm{R} (u, \lambda) \circ \sum_ {k \in \mathbf {N}} (\lambda - \mu) ^ {k} \mathrm{R} (u, \lambda) ^ {k},
\]

which implies that the resolvent of u is holomorphic in U.

The assertion \(b\) ) follows immediately from \(a\) ). If \(\mathrm{Sp}(u)\) is nonempty, the distance from \(\lambda\) to \(\mathrm{Sp}(u)\) is strictly positive (TG, IX, p. 13, prop. 2), whence \(c\) ).

The last part of the assertion \(d\) ) is proved by induction on \(k\) , the case \(k = 1\) being a consequence of formula (8), p. 245.

PROPOSITION 15. --- Let u be a closed operator with dense domain on a complex Banach space E, and let λ be a complex number belonging to the resolvent set of u.

\begin{enumerate}
\def\labelenumi{\alph{enumi})}
\item
  The subset \(\mathrm{Sp}(\mathrm{R}(u,\lambda)) \setminus \{0\}\) of \(\mathbf{C}\) is the image of the spectrum of \(u\) by the map \(\mu \mapsto (\lambda -\mu)^{-1}\) from \(\mathbf{C} \setminus \{\lambda \}\) into \(\mathbf{C}^*\) ;
\item
  For every \(\mu \neq \lambda\) in \(\mathbf{C}\) , one has
\end{enumerate}

\[
\operatorname{Ker} \left(\mu 1 _ {\mathrm{E}} - u\right) = \operatorname{Ker} \left(\left(\lambda - \mu\right) ^ {- 1} 1 _ {\mathrm{E}} - \mathrm{R} (u, \lambda)\right).
\]

Let us prove assertion \(a\) ). For every \(\mu \neq \lambda\) , we compute

\[
\mu 1 _ {\mathrm{E}} - u = (\lambda - \mu) ((\lambda - \mu) ^ {- 1} 1 _ {\mathrm{E}} - \mathrm{R} (u, \lambda)) (\lambda 1 _ {\mathrm{E}} - u).
\]

Since \(\lambda \notin \mathrm{Sp}(u)\) and \(\mu \neq \lambda\) , the linear map \((\lambda - \mu)(\lambda 1_{\mathrm{E}} - u)\) is a bijection from \(\mathrm{dom}(u)\) onto E. Consequently, this formula implies that \(\mu 1_{\mathrm{E}} - u\) is a bijection from \(\mathrm{dom}(u)\) onto E if and only if \((\lambda - \mu)^{-1} 1_{\mathrm{E}} - \mathrm{R}(u, \lambda)\) is a bijection from E onto E, which implies the assertion.

Let us prove \(b\) ). If \(\mu \neq \lambda\) and \(x \in \mathrm{Ker}((\lambda - \mu)^{-1}1_{\mathrm{E}} - \mathrm{R}(u, \lambda))\) , we have \(x \in \mathrm{dom}(u)\) and the formula \(1_{\mathrm{E}} = (\lambda 1_{\mathrm{E}} - u) \circ \mathrm{R}(u, \lambda)\) implies that \(x \in \mathrm{Ker}(\mu 1_{\mathrm{E}} - u)\) . Conversely, if \(x \in \mathrm{Ker}(\mu 1_{\mathrm{E}} - u)\) and \(\mu \neq \lambda\) , the formula \(1_{\mathrm{dom}(u)} = \mathrm{R}(u, \lambda) \circ (\lambda 1_{\mathrm{E}} - u)\) implies \(\mathrm{R}(u, \lambda)(x) = (\lambda - \mu)^{-1}x\) .

PROPOSITION 16. --- Let u be a closed operator with dense domain on a complex Hilbert space E. The spectrum of \(u^{*}\) is the image of the spectrum of u under complex conjugation and, for every element \(\lambda\) of the resolvent set of u, we have \(\mathrm{R}(u,\lambda)^{*}=\mathrm{R}(u^{*},\overline{\lambda})\) . In particular, if u is self-adjoint, the endomorphism \(\mathrm{R}(u,\lambda)\) is normal for every \(\lambda\notin\mathrm{Sp}(u)\) .

Let \(\lambda\in\mathbf{C} - \mathrm{Sp}(u)\) be an element of the resolvent set of \(u\) . We have \(u=\lambda1_{\mathrm{E}}-\mathrm{R}(u,\lambda)^{-1}\) , hence

\[
u ^ {*} = \overline {{\lambda}} 1 _ {\mathrm{E}} - (\mathrm{R} (u, \lambda) ^ {- 1}) ^ {*} = \overline {{\lambda}} 1 _ {\mathrm{E}} - (\mathrm{R} (u, \lambda) ^ {*}) ^ {- 1}
\]

(IV, p. 236 and prop. 9 of IV, p. 239). By Remark 5, it follows that \(\overline{\lambda} \in \mathbf{C} - \mathrm{Sp}(u^{*})\) and that \(\mathrm{R}(u, \lambda)^{*} = \mathrm{R}(u^{*}, \overline{\lambda})\) . Consequently, the spectrum of \(u^{*}\) is contained in the image of \(\mathrm{Sp}(u)\) by complex conjugation. The equality is obtained by applying this property to \(u^{*}\) since \(u^{**} = u\) (corollary of Proposition 8 in IV, p. 237). The last assertion then follows from formula (9), p. 245.

COROLLARY. --- Let u be a partial self-adjoint operator on a complex Hilbert space E. If E is not zero, then the spectrum of u is not empty.

Suppose that \(\operatorname{Sp}(u)\) is empty. Then u is injective and \(u^{-1} = -\operatorname{R}(u,0)\) is an injective endomorphism of E such that \(\operatorname{Sp}(u^{-1}) \subset \{0\}\) (prop. 15, a)), hence \(\operatorname{Sp}(u^{-1}) = \{0\}\) (I, p. 26, cor. 1). Since \(u^{-1}\) is normal (prop. 16), this implies that \(u^{-1}\) is zero (I, p. 110, example 1), which is a contradiction.

Lemma 5. --- Let E be a complex Hilbert space and let u be a closed densely defined partial operator on E. Let \(\lambda \in \mathbf{C}\) . Suppose there exists a real number \(c > 0\) such that

\[
\| u (x) - \lambda x \| \geqslant c \| x \| \quad \text{for all } x \in \operatorname{dom} (u),\tag{10}
\]

\[
\| u ^ {*} (x) - \overline {{\lambda}} x \| \geqslant c \| x \| \quad \text{for all } x \in \operatorname{dom} (u ^ {*}),\tag{11}
\]

Then \(\lambda\) belongs to the resolvent set of \(u\) and \(\| \mathrm{R}(u,\lambda)\| \leqslant c^{-1}\) .

The hypothesis implies that \(u - \lambda 1_{\mathrm{E}}\) and \(u^{*} - \overline{\lambda} 1_{\mathrm{E}}\) are injective. Let F be the image of \(u - \lambda 1_{\mathrm{E}}\) . The space F is dense in E, since its orthogonal is equal to \(\operatorname{Ker}(u^{*} - \overline{\lambda} 1_{\mathrm{E}})\) (prop. 7, c) of IV, p. 236), which is zero.

Let us prove that the space F is closed. Let \((x_{n})_{n\in\mathbf{N}}\) be a sequence in dom(u) such that the sequence \((u(x_{n})-\lambda x_{n})_{n\in\mathbf{N}}\) converges to \(y\in F\) Inequality (10) implies that the sequence \((x_{n})_{n\in\mathbf{N}}\) is a Cauchy sequence in E. Let \(x\in E\) be its limit. The sequence \((x_{n},u(x_{n}))\) converges to \((x,y+\lambda x)\) in E×E; since the graph of u is closed, we therefore have \(x\in\text{dom}(u)\) and \(u(x)=y+\lambda x\) , which proves that \(y\in F\) .

We conclude that F = E. Thus, the partial operator \(u - \lambda1_{E}\) is bijective, whence \(\lambda \notin \operatorname{Sp}(u)\) (Remark 3). Inequality (10) then implies that \(\|R(u, \lambda)\| \leqslant c^{-1}\) .

PROPOSITION 17. — Let E be a complex Hilbert space and let u be a self-adjoint partial operator on E.

\begin{enumerate}
\def\labelenumi{\alph{enumi})}
\item
  The spectrum of u is contained in \(\mathbf{R}\) ;
\item
  If u is positive, then the spectrum of u is contained in \(\mathbf{R}_{+}\) ;
\item
  Suppose that E is nonzero. Let \(\lambda \notin \operatorname{Sp}(u)\) and let \(\delta > 0\) be the distance from \(\lambda\) to the spectrum of \(u\) . We have \(\| \mathrm{R}(u, \lambda) \| = \delta^{-1}\) .
\end{enumerate}

Let \((a,b)\in\mathbf{R}\times\mathbf{R}\) and \(\lambda=a+ib\) . Let \(x\in\mathrm{dom}(u)\) . Since u is self-adjoint, we have \(\langle x|u(x)\rangle\in\mathbf{R}\) , whence

\[
\begin{array}{r} \| u (x) - \lambda x \| ^ {2} = \| u (x) \| ^ {2} - 2 a \langle x | u (x) \rangle + (a ^ {2} + b ^ {2}) \| x \| ^ {2} \\ = \| u (x) - \overline {{\lambda}} x \| ^ {2}. \end{array}
\]

Suppose that \(b \neq 0\) . We then obtain

\[
\begin{array}{c} \| u (x) - \lambda x \| ^ {2} = \| u (x) - \overline {{\lambda}} x \| ^ {2} \geqslant (\| u (x) \| - a \| x \|) ^ {2} + b ^ {2} \| x \| ^ {2} \\ \geqslant b ^ {2} \| x \| ^ {2}. \end{array}
\]

By Lemma 5, we therefore have \(\lambda \notin \operatorname{Sp}(u)\) , whence assertion \(a\) ).

Suppose that u is also positive. If b = 0 and a \textless{} 0, we obtain similarly for \(x \in \text{dom}(u)\) the inequality

\[
\left\| u (x) - \lambda x \right\| ^ {2} = \left\| u (x) - \overline {{\lambda}} x \right\| ^ {2} \geqslant \left(\left\| u (x) \right\| - a \| x \|\right) ^ {2} \geqslant a ^ {2} \| x \| ^ {2},
\]

hence \(\lambda\notin\mathrm{Sp}(u)\) (loc. cit.), which proves \(b\) ).

Let us finally prove \(c\) ). By prop. 16, the resolvent \(\mathrm{R}(u,\lambda)\) is a normal endomorphism of E. Its norm is therefore equal to its spectral radius (cor. 1 of I, p. 108), whence

\[
\| \mathrm{R} (u, \lambda) \| = \sup _ {\mu \in \operatorname{Sp} (\mathrm{R} (u, \lambda))} | \mu |
\]

(theorem 1 of I, p. 24). The spectrum of R(u,λ) cannot be reduced to \{0\}, for in that case one would have \(\| \mathrm{R}(u,\lambda)\| = 0\), so \(\operatorname{dom}(u)\), the image of R(u,λ), would be zero, and E as well. Prop. 15 therefore implies that

\[
\| \mathrm{R} (u, \lambda) \| = \sup _ {\mu \in \mathrm{Sp} (u)} \frac {1}{| \lambda - \mu |} = \frac {1}{\delta}.
\]

COROLLARY. --- Let u be a self-adjoint operator on E.

\begin{enumerate}
\def\labelenumi{\alph{enumi})}
\item
  The residual spectrum of u is empty;
\item
  For all \(\lambda \neq \mu\) in \(\mathbf{C}\) , the eigenspaces of \(u\) relative to \(\lambda\) and \(\mu\) are orthogonal.
\end{enumerate}

Let us prove a). Let \(\lambda\) belong to the spectrum of \(u\) ; it is a real number (prop. 17). We have \(\mathrm{Ker}(\lambda 1_{\mathrm{E}} - u) = \mathrm{Im}(\lambda 1_{\mathrm{E}} - u)^{\circ}\) (prop. 7 of IV, p. 236), hence \(\lambda\) is not an eigenvalue of \(u\) if the partial operator \(\lambda 1_{\mathrm{E}} - u\) has dense image. This implies by definition that the residual spectrum of \(u\) is empty.

Let us prove \(b\) ). By the proposition, we may assume that \(\lambda\) and \(\mu\) are real. Let \(x \in \operatorname{dom}(u)\) such that \(u(x) = \lambda x\) and \(y \in \operatorname{dom}(u)\) such that \(u(y) = \mu y\) . We then have

\[
\lambda \langle x \mid y \rangle = \langle u (x) \mid y \rangle = \langle x \mid u (y) \rangle = \mu \langle x \mid y \rangle ,
\]

whence \(\langle x \mid y \rangle = 0.\)

Remark. --- If u is a closed symmetric operator that is not self-adjoint, its spectrum is not contained in R (cf. Cor. 10 of IV, p. 257 below), and it is possible that the eigenspaces of u corresponding to λ and \(\overline{\lambda}\) need not be orthogonal (exercise 11 of IV, p. 347).

\subsection{Pseudospectrum}

DEFINITION 9. --- Let E be a complex Banach space and let u be a closed partial operator on E. Let \(\varepsilon\) be a strictly positive real number. We call the \(\varepsilon\)-pseudospectrum of u the union of the spectrum of u and the set of complex numbers \(\lambda\) belonging to the resolvent set of u such that \(\| \mathrm{R}(u,\lambda)\| >\varepsilon^{-1}\) . We denote this set by \(\mathrm{PSp}_{\varepsilon}(u)\) .

\pandocbounded{\includegraphics[keepaspectratio]{images/part002_14758b8b.jpg}}

Some authors define the \(\varepsilon\) -pseudo-spectrum of \(u\) as the set \(\mathrm{T}_{\varepsilon}(u)\), which is the union of \(\mathrm{Sp}(u)\) and the set of \(\lambda \in \mathbf{C} - \mathrm{Sp}(u)\) such that \(\| \mathrm{R}(u,\lambda)\| \geqslant \varepsilon^{-1}\) . The closure of \(\mathrm{PSp}_{\varepsilon}(u)\) is contained in \(\mathrm{T}_{\varepsilon}(u)\) but this inclusion can be strict, even if E is a Hilbert space (cf. exercises 18 of IV, p. 348 and 19 of IV, p. 349).

PROPOSITION 18. --- Let E be a complex Banach space and u a closed partial operator on E. Let \(\varepsilon \in \mathbf{R}_{+}^{*}\) . The set \(\mathrm{PSp}_{\varepsilon}(u)\) is an open subset of C. It is nonempty if E is nonzero.

If E is zero, then \(\mathrm{PSp}_{\varepsilon}(u)\) is empty. Suppose that E is nonzero. By prop. 14 of IV, p. 246, the set \(\mathrm{PSp}_{\varepsilon}(u)\) is an open subset of C.

The set \(\mathrm{PSp}_{\varepsilon}(u)\) is not empty if the spectrum of \(u\) is not empty. If \(\mathrm{Sp}(u)\) is empty, then Liouville's theorem (VAR, R1, p. 29, 3.3.6) implies that the holomorphic function on \(\mathbf{C}\) defined by \(\lambda \mapsto \mathrm{R}(u,\lambda)\) is not bounded, so there exists \(\lambda\) in \(\mathbf{C}\) such that \(\| \mathrm{R}(u,\lambda)\| >\varepsilon^{-1}\) .

PROPOSITION 19. --- Let E be a complex Banach space and u a closed partial operator on E. Then we have \(\mathrm{PSp}_{\varepsilon}(u) \subset \mathrm{PSp}_{\delta}(u)\) if \(0 < \varepsilon < \delta\) and

\[
\bigcap_ {\varepsilon \in \mathbf {R} _ {+} ^ {*}} \mathrm{PSp} _ {\varepsilon} (u) = \mathrm{Sp} (u).
\]

The first assertion follows from the definition. For the second, the spectrum of u is contained in \(\mathrm{PSp}_{\varepsilon}(u)\) for all \(\varepsilon > 0\) by definition, and if \(\lambda \notin \mathrm{Sp}(u)\) , then \(\lambda \notin \mathrm{PSp}_{\varepsilon}(u)\) when \(\varepsilon < \|R(u, \lambda)\|^{-1}\) .

PROPOSITION 20. --- Let E be a complex Banach space and u a closed partial operator on E. Let \(\varepsilon\in\mathbf{R}_{+}^{*}\) . For every \(\lambda\in\mathbf{C}\) the following conditions are equivalent:

\begin{enumerate}
\def\labelenumi{(\roman{enumi})}
\item
  One has \(\lambda \in \mathrm{PSp}_{\varepsilon}(u)\) ;
\item
  Either \(\lambda \in \mathrm{Sp}(u)\), or there exists \(x \in \mathrm{dom}(u)\) such that \(\| x \| = 1\) and \(\| (\lambda 1_{\mathrm{E}} - u)(x) \| < \varepsilon\) ;
\item
  There exists \(v \in \mathcal{L}(\mathrm{E})\) such that \(\| v \| < \varepsilon\) and \(\lambda \in \mathrm{Sp}(u + v)\) .
\end{enumerate}

We may assume that E is nonzero. Condition (ii) is a reformulation of the definition, and hence of condition (i).

Suppose condition (i) holds, and let us prove (iii). If \(\lambda\) belongs to the spectrum of \(u\) one may take \(v = 0\) in (iii).

Suppose therefore that \(\lambda \notin \mathrm{Sp}(u)\) . By definition of \(\mathrm{PSp}_{\varepsilon}(u)\) , there exists \(y \in \mathrm{E}\) such that \(\| y \| = 1\) and \(\| \mathrm{R}(u, \lambda)y \| > \varepsilon^{-1}\) . Set \(x = \mathrm{R}(u, \lambda)y\) . We have \(x \neq 0\) . By the Hahn-Banach theorem (EVT, II, p. 67, cor. 2), the linear form \(\ell\) on \(\mathbf{C}x\) such that \(\ell(x) = 1\) admits a continuous extension \(\ell_1 \in \mathrm{E}'\) such that \(\| \ell_1 \| = \| \ell \|\) ; we therefore have \(\| \ell_1 \| = \| x \|^{-1} < \varepsilon\) . For every \(e \in \mathrm{E}\) set \(v(e) = \ell_1(e)y\) . We have \(v \in \mathcal{L}(\mathrm{E})\) and \(v(x) = y\) . We obtain \((u + v)x = u(x) + y = \lambda x\) , hence \(\lambda \in \mathrm{Sp}(u + v)\) . Since moreover \(\| v \| = \| \ell_1 \| < \varepsilon\) the condition (iii) is satisfied.

Conversely, let \(v \in \mathcal{L}(\mathrm{E})\) such that \(\|v\| < \varepsilon\) and \(\lambda \in \mathrm{Sp}(u + v)\) . The partial operator \(\lambda1_{\mathrm{E}} - (u + v)\) is therefore either not injective or does not have a continuous inverse; by Lemma 4 of IV, p. 243, applied to the injective partial operator \(\lambda1_{E} - u\) and to \(-v\), we therefore have \(\|v\| \geqslant \|R(u, \lambda)\|^{-1}\) . It follows that condition (iii) implies (i).

COROLLARY. — Let E be a complex Banach space and u a closed partial operator on E. Let \(\varepsilon > 0\) .

\begin{enumerate}
\def\labelenumi{\alph{enumi})}
\item
  For every \(v \in \mathcal{L}(\mathrm{E})\) , we have \(\mathrm{PSp}_{\varepsilon}(u) \subset \mathrm{PSp}_{\varepsilon + \| v\|}(u + v)\) ;
\item
  Let \(\delta > 0\) and let \(D_{\delta}\) be the open disk with center 0 and radius \(\delta\) in \(\mathbf{C}\) . We have \(\mathrm{PSp}_{\varepsilon}(u) + D_{\delta} \subset \mathrm{PSp}_{\varepsilon + \delta}(u)\) .
\end{enumerate}

Let \(\lambda \in \mathrm{PSp}_{\varepsilon}(u)\) . There exists an endomorphism \(w \in \mathcal{L}(\mathrm{E})\) such that \(\|w\| < \varepsilon\) and \(\lambda \in \mathrm{Sp}(u+w)\) (prop. 20). Since \(u+w = (u+v)+(w-v)\) and \(\|w-v\| < \varepsilon + \|v\|\) , we have \(\lambda \in \mathrm{PSp}_{\varepsilon+ \|v\|}(u+v)\) (loc. cit.).

Let \(\mu \in \mathrm{D}_{\delta}\) ; we have \(\lambda + \mu \in \mathrm{Sp}(u + (w + \mu 1_{\mathrm{E}}))\) and \(\| w + \mu 1_{\mathrm{E}} \| < \varepsilon + \delta\) , hence \(\lambda + \mu \in \mathrm{PSp}_{\varepsilon + \delta}(u)\) (loc. cit.).

PROPOSITION 21. --- Let E be a complex Banach space and u a closed partial operator on E. Let \(\varepsilon\in\mathbf{R}_{+}^{*}\) . Every bounded connected component of \(\mathrm{PSp}_{\varepsilon}(u)\) meets the spectrum of \(u\) .

Let U be a connected component of \(\mathrm{PSp}_{\varepsilon}(u)\) which does not meet \(\mathrm{Sp}(u)\) . The set U is open and closed in \(\mathrm{PSp}_{\varepsilon}(u)\) and its closure \(\overline{U}\) in C therefore satisfies \(\overline{U} \cap \mathrm{PSp}_{\varepsilon}(u) = U\) .

Since \(\mathrm{Sp}(u)\) is contained in \(\mathrm{PSp}_{\varepsilon}(u)\) and U does not meet \(\mathrm{Sp}(u)\) , this equality proves that the set \(\overline{U}\) is disjoint from \(\mathrm{Sp}(u)\) , hence contained in the resolvent set of u. Moreover, it implies that the set \(\overline{U}-U\) does not meet \(\mathrm{PSp}_{\varepsilon}(u)\) .

We therefore have \(\|\mathrm{R}(u,\lambda)\| \leqslant \varepsilon^{-1}\) for all \(\lambda\) in \(\overline{U} - U\) , while \(\|\mathrm{R}(u,\lambda)\| > \varepsilon^{-1}\) for \(\lambda \in U\) . If the set U is bounded, its closure \(\overline{U}\) is compact and there exists \(\lambda_{0} \in \overline{U}\) such that \(\|\mathrm{R}(u,\lambda)\| \leqslant \|\mathrm{R}(u,\lambda_{0})\|\) for \(\lambda \in \overline{U}\) . The foregoing implies that \(\lambda_{0} \in U\) which contradicts the maximum principle (VAR, R1, p. 29, 3.3.7) since the resolvent of u is holomorphic in the resolvent set of u (prop. 14 of IV, p. 246).

\subsection{Multiplication operators}

Let X be a locally compact space and \(\mu\) a positive measure on X. We consider the multiplication operators on \(\mathrm{L}^{2}(\mathrm{X},\mu)\) ; these are closed operators with dense domain (prop. 5 of IV, p. 232). For any \(\mu\)-measurable function g on X, we will denote by \(m_{g}\) the partial multiplication operator by g on \(\mathrm{L}^{2}(\mathrm{X},\mu)\) .

PROPOSITION 22. — Let \(g\) be a function \(\mu\)-measurable on X.

\begin{enumerate}
\def\labelenumi{\alph{enumi})}
\item
  The spectrum of \(m_{g}\) is the \(\mu\)-essential image S of \(g\) ;
\item
  Let \(\lambda \in \mathbf{C} - \mathrm{Sp}(m_g)\) . The resolvent \(\mathrm{R}(m_g, \lambda)\) is the multiplication operator \(m_h\) , where \(h\) is the function on X defined by \(h(x) = 0\) if \(g(x) = \lambda\) and \(h(x) = (\lambda - g(x))^{-1}\) otherwise.
\end{enumerate}

Let us prove that \(\mathbf{C} - \mathbf{S}\) is contained in the resolvent set of \(m_g\) . Let \(\lambda \in \mathbf{C} - \mathbf{S}\) . There exists an open neighborhood U of \(\lambda\) such that the subset \(\mathrm{Y} = g^{-1}(\mathrm{U})\) of X is locally \(\mu\) -negligible. The function \(k\) defined on X by \(k(x) = (\lambda - g(x))^{-1}\) if \(x \notin \mathrm{Y}\) and \(k(x) = 0\) if \(x \in \mathrm{Y}\) belongs to \(\mathcal{L}^{\infty}(\mathrm{X}, \mu)\) (Lemma 5 of IV, p. 184); the operator of multiplication by \(k\) is therefore an endomorphism of \(\mathrm{L}^{2}(\mathrm{X}, \mu)\) .

Since \(|gk| \leqslant 1 + |\lambda k|\) , we have

\[
| g k f | \leqslant | f | + | \lambda k f |
\]

for \(f \in \mathcal{L}^2(\mathrm{X},\mu)\) , which implies that the image of \(m_k\) is contained in the domain of \(m_g\) . Conversely, let \(f \in \mathcal{L}^2(\mathrm{X},\mu)\) whose class \(\widetilde{f}\) belongs to the domain of \(m_g\) . Outside the locally \(\mu\) -negligible set Y, we have \(f(x) = k(x)(\lambda - g(x))f(x)\) , hence \(\widetilde{f}\) is in the image of \(m_k\) . The same formula proves that \(\lambda\) belongs to the resolvent set of \(m_g\) and that \(\mathrm{R}(m_g,\lambda) = m_k\) . Since the set Y is locally \(\mu\) -negligible, the multiplication operator \(m_k\) coincides with the operator \(m_h\) described in assertion \(b\) ).

Let us prove conversely that \(\mathbf{C} - \mathrm{Sp}(m_g)\) is contained in \(\mathbf{C} - \mathrm{S}\) . Let \(\lambda \in \mathbf{C} - \mathrm{Sp}(m_g)\) . Let \(\mathrm{M} > \| \mathrm{R}(m_g, \lambda) \|\) be a real number. Let Y denote the set of \(x \in \mathrm{X}\) such that \(|\lambda - g(x)| < \mathrm{M}^{-1}\) . Let us prove that Y is locally \(\mu\) -negligible, which will imply that \(\lambda\) does not belong to S, and will conclude the proof.

Let \(K\) be a compact subset of \(X\). Let \(\varphi\) be the characteristic function of \(Y \cap K\); this is an element of \(\mathcal{L}^{2}(X,\mu)\), and we denote its class in \(\mathrm{L}^{2}(X,\mu)\) by \(\widetilde{\varphi}\). Let \(\psi\) be a function in \(\mathcal{L}^{2}(X,\mu)\) whose class in \(\mathrm{L}^{2}(X,\mu)\) is \(\mathrm{R}(m_g,\lambda)(\widetilde{\varphi})\). We have \(\mathrm{R}(m_g,\lambda)(\widetilde{\varphi}) \in \operatorname{dom}(m_g)\) and \((\lambda - m_g)(\mathrm{R}(m_g,\lambda)(\widetilde{\varphi})) = \widetilde{\varphi}\), hence \((\lambda - g(x))\psi(x) = 1\) for \(\mu\)-almost all \(x \in Y \cap K\). This implies

\[
\| \mathrm{R} (m _ {g}, \lambda) \| ^ {2} \| \widetilde {\varphi} \| ^ {2} \geqslant \| \mathrm{R} (m _ {g}, \lambda) (\widetilde {\varphi}) \| ^ {2} \geqslant \mathrm{M} ^ {2} \mu (\mathrm{Y} \cap \mathrm{K}) = \mathrm{M} ^ {2} \| \widetilde {\varphi} \| ^ {2}.
\]

In view of the choice of M, this means that \(\varphi\) is zero \(\mu\) -almost everywhere. Thus, the set Y \(\cap\) K is \(\mu\) -negligible. The set Y is therefore locally \(\mu\) -negligible (INT, IV, p. 172, § 5, n° 2, prop. 5).

PROPOSITION 23. --- Let g be a μ-measurable function on X. The adjoint of the multiplication operator \(m_{g}\) is \(m_{\overline{g}}\) .

For every integer \(n \geqslant 1\) , let \(\varphi_{n}\) be the characteristic function of the set of elements \(x \in X\) such that \(|g(x)| \leqslant n\) and let \(\widetilde{\varphi}_{n}\) be its class in \(\mathrm{L}^{2}(\mathrm{X}, \mu)\) . Let \(f \in \mathcal{L}^{2}(\mathrm{X}, \mu)\) whose class \(\widetilde{f}\) belongs to \(\operatorname{dom}(m_{g}^{*})\) and let \(\psi\) be a function whose class is \(m_{g}^{*}(\widetilde{f})\) .

For every \(h \in \mathcal{L}^2(\mathrm{X}, \mu)\) whose class \(\widetilde{h}\) belongs to \(\operatorname{dom}(m_g)\) , we also have \(\widetilde{h}\widetilde{\varphi}_n \in \operatorname{dom}(m_g)\) , and hence \(\langle \widetilde{f} | m_g(\widetilde{h}\widetilde{\varphi}_n) \rangle = \langle m_g^*(\widetilde{f}) | \widetilde{h}\widetilde{\varphi}_n \rangle\) . This entails the equality

\[
\int_ {\mathrm{X}} \overline {{(f \overline {{g}} - \psi)}} \varphi_ {n} h d \mu = 0.
\]

Since the domain of \(m_{g}\) is dense in \(\mathrm{L}^{2}(\mathrm{X},\mu)\) , we deduce that \((f\overline{g}-\psi)\varphi_{n}\) is zero \(\mu\) -almost everywhere. Since n is arbitrary, this means that \(m_{g}^{*}(\widetilde{f})\) is the class in \(\mathrm{L}^{2}(\mathrm{X},\mu)\) of \(f\overline{g}\) . In particular, since \(m_{g}^{*}(\widetilde{f}) \in \mathrm{L}^{2}(\mathrm{X},\mu)\) we conclude that f belongs to the domain of \(m_{\overline{g}}\) and that \(m_{g}^{*}(\widetilde{f}) = m_{\overline{g}}(\widetilde{f})\) .

The adjoint of \(m_{g}\) is therefore an extension of \(m_{\overline{g}}\) . Moreover, we have

\[
\langle f \mid m _ {g} (h) \rangle = \int_ {\mathrm{X}} \overline {{f}} \cdot (g h) d \mu = \int_ {\mathrm{X}} f \overline {{g}} h d \mu
\]

for all \(f \in \mathrm{L}^{2}(\mathrm{X}, \mu)\) and \(h \in \mathrm{dom}(m_{g})\) , which proves that the linear form \(h \mapsto \langle f | m_{g}(h) \rangle\) is continuous when \(m_{\overline{g}}(f) g\) belongs to \(\mathrm{L}^{2}(\mathrm{X}, \mu)\) . Consequently, the domain of \(m_{\overline{g}}\) is contained in that of \(m_{g}^{*}\) , which concludes the proof.

COROLLARY. --- Let g be a μ-measurable function on X. The multiplication operator \(m_{g}\) on \(\mathrm{L}^{2}(\mathrm{X},\mu)\) is self-adjoint (resp. positive) if and only if the function g is locally μ-almost everywhere real-valued (resp. locally μ-almost everywhere positive).

The first assertion follows from the proposition. If \(g\) is locally \(\mu\)-almost everywhere positive, we have \(\langle f|m_{g}(f)\rangle=\int_{\mathrm{X}}g|f|^{2}d\mu\geqslant0\) for all \(f\in\mathrm{L}^{2}(\mathrm{X},\mu)\) , so the partial operator \(m_{g}\) is positive.

Conversely, if \(m_{g}\) is positive, then its spectrum is contained in \(\mathbf{R}_{+}\) (prop. 17 of IV, p. 248); since it is the \(\mu\)-essential image of g (prop. 22, a)), this means that g is locally \(\mu\) -almost everywhere positive.

Lemma 6. --- Let \(g_{1}\) and \(g_{2}\) be functions \(\mu\)-measurable on X.

\begin{enumerate}
\def\labelenumi{\alph{enumi})}
\item
  The partial operator \(m_{g_{1}} + m_{g_{2}}\) is closable and its closure is the multiplication operator \(m_{g_{1}+g_{2}}\) ;
\item
  One has \(m_{g_{1}} \circ m_{g_{2}} \subset m_{g_{1}g_{2}}\) ;
\item
  Suppose that \(g_{1}\) is bounded. We then have \(m_{g_{2}} \circ m_{g_{1}} = m_{g_{1}g_{2}}\) . Moreover, the domain of \(m_{g_{2}}\) is contained in the domain of \(m_{g_{1}g_{2}}\) , and \(m_{g_{1}} \circ m_{g_{2}}\) is the reduction of \(m_{g_{1}g_{2}}\) to \(\text{dom}(m_{g_{2}})\) .
\end{enumerate}

It is elementary that \(m_{g_{1}+g_{2}}\) is an extension of \(m_{g_{1}} + m_{g_{2}}\) ; this latter operator is therefore closable, and \(\overline{m_{g_{1}} + m_{g_{2}}} \subset m_{g_{1}+g_{2}}\) .

Let \(f \in \mathcal{L}^2(\mathrm{X},\mu)\) such that the function \(h = (g_1 + g_2)f\) belongs to \(\mathcal{L}^2(\mathrm{X},\mu)\) . For every integer \(n \geqslant 1\) , denote by \(\mathrm{X}_n\) the set of \(x \in \mathrm{X}\) such that \(|g_1(x)| + |g_2(x)| \leqslant n\) and let \(\varphi_n\) be the characteristic function of \(\mathrm{X}_n\) . We have \(\varphi_n f \in \mathrm{dom}(m_{g_1} + m_{g_2})\) . Since \((\varphi_n f)(x)\) converges to \(f(x)\) for all \(x \in \mathrm{X}\) and since for every \(n \in \mathbf{N}\) we have \(|\varphi_n f| \leqslant |f|\) and \(|(g_1 + g_2)\varphi_n f| \leqslant |(g_1 + g_2)f| = |h|\) , with \(h \in \mathcal{L}^2(\mathrm{X},\mu)\) by hypothesis, Lebesgue's theorem (INT, IV, p. 137, § 3, n° 7, th. 6) implies that the sequence of pairs of classes in \(\mathrm{L}^{2}(\mathrm{X},\mu)\times\mathrm{L}^{2}(\mathrm{X},\mu)\) of \((\varphi_{n}f,(g_{1}+g_{2})\varphi_{n}f)\) which belong to the graph of \(m_{g_{1}}+m_{g_{2}}\) converges to the pair of classes of \((f,h)\) in \(\mathrm{L}^{2}(\mathrm{X},\mu)\) . The closure of \(m_{g_{1}}+m_{g_{2}}\) is therefore indeed equal to \(m_{g_{1}+g_{2}}\) .

It is elementary that \(m_{g_{1}} \circ m_{g_{2}} \subset m_{g_{1}g_{2}}\) . Suppose that \(g_{1}\) is bounded, so that \(m_{g_{1}}\) is a continuous linear map on \(\mathrm{L}^{2}(\mathrm{X},\mu)\) . Then the domain of \(m_{g_{2}} \circ m_{g_{1}}\) is the set of equivalence classes of functions \(f \in \mathcal{L}^{2}(\mathrm{X},\mu)\) such that \(g_{2}(g_{1}f) \in \mathcal{L}^{2}(\mathrm{X},\mu)\) , that is, the domain of \(m_{g_{1}g_{2}}\) . We therefore have \(m_{g_{2}} \circ m_{g_{1}} = m_{g_{1}g_{2}}\) .

It is also elementary that \(\operatorname{dom}(m_{g_{1}} \circ m_{g_{2}}) = \operatorname{dom}(m_{g_{2}})\) is contained in \(\operatorname{dom}(m_{g_{1}g_{2}})\) , and that the restriction of \(m_{g_{1}g_{2}}\) to this space is equal to \(m_{g_{1}} \circ m_{g_{2}}\) .

One should beware of believing that \(m_{g_{1}} \circ m_{g_{2}} = m_{g_{1}g_{2}}\) in general (exercise 10 of IV, p. 347). Nevertheless, we have the following partial result:

PROPOSITION 24. --- Let g be a μ-measurable function on X.

\begin{enumerate}
\def\labelenumi{\alph{enumi})}
\item
  One has \(m_{\overline{g}} \circ m_g = m_{|g|^2}\) ;
\item
  For all integers \(k, \ell \in \mathbf{N}\) , we have \(m_{g^k\overline{g}^\ell} = m_g^k m_{\overline{g}}^\ell\) .
\end{enumerate}

One has \(m_{\overline{g}} \circ m_g \subset m_{|g|^2}\) by Lemma 6. Conversely, from the inequality \(|g| \leqslant 1 + |g|^2\) we deduce that \(\text{dom}(m_{|g|^2}) = \text{dom}(m_{\overline{g}} \circ m_g)\) , whence the first assertion.

Let \(k, \ell \in \mathbf{N}\) . We have \(m_g^k m_{\overline{g}}^\ell \subset m_{g^k \overline{g}^\ell}\) (loc. cit.). The domain of \(m_g^k m_{\overline{g}}^\ell\) is the set of classes in \(\mathrm{L}^2(\mathrm{X}, \mu)\) of functions \(h \in \mathcal{L}^2(\mathrm{X}, \mu)\) such that \(|g|^j h\) belongs to \(\mathcal{L}^2(\mathrm{X}, \mu)\) for every integer \(j\) such that \(0 \leqslant j \leqslant k + \ell\) . The inequalities

\[
\left| g ^ {j} h \right| \leqslant | h | + \left| g ^ {k + \ell} h \right|,
\]

valid for \(0 \leqslant j \leqslant k + \ell\) , allow one to observe that \(\mathrm{dom}(m_g^k m_{\overline{g}}^\ell)\) is equal to \(\mathrm{dom}(m_{g^k\overline{g}^\ell})\) , whence assertion \(b\) ).

\subsection{Self-adjoint extensions of a symmetric operator}

In this section, we propose to classify the self-adjoint extensions of symmetric operators on a complex Hilbert space, and in particular to find conditions ensuring the existence of a self-adjoint extension of such an operator.

Let E be a complex Hilbert space. For every partial operator u on E, we shall write in this section \(u + i\) and \(u - i\) instead of \(u + i1_{E}\) and \(u - i1_{E}\) .

Let u be a closed operator with dense domain on E. We recall (cf. definition 4 of IV, p. 230 and proposition 4 of IV, p. 230) that we denote by \(E_{u}\) the Hilbert space dom(u) equipped with the Hermitian form

\[
(x \mid y) _ {u} = \langle x \mid y \rangle + \langle u (x) \mid u (y) \rangle .
\]

We denote \(\|x\|_{u}\) the norm of an element x of the Hilbert space \(E_{u}\) . Every subspace of dom(u) that is closed in E is closed in \(E_{u}\) (IV, p. 230, remark 1).

DEFINITION 10. --- Let u be a closed symmetric operator on a complex Hilbert space E. Let \(\mathrm{E}_{+} = \mathrm{Ker}(u^{*} - i)\) and \(\mathrm{E}_{-} = \mathrm{Ker}(u^{*} + i)\) . The pair \((\mathrm{E}_{+}, \mathrm{E}_{-})\) of subspaces of \(\mathrm{dom}(u^{*})\) is called the deficiency pair of u. The subspaces \(E_{+}\) and \(E_{-}\) are closed subspaces of E (prop. 3 of IV, p. 229). The pair \((\dim(\mathrm{E}_{+}), \dim(\mathrm{E}_{-}))\) of Hilbertian dimensions of these subspaces is called the deficiency index of u.

PROPOSITION 25. --- Let u be a closed symmetric operator on a complex Hilbert space E. We have

\[
\begin{array}{l} \operatorname{Ker} (u ^ {*} - i) ^ {\circ} = \operatorname{Im} (u + i), \\ \operatorname{Ker} (u ^ {*} + i) ^ {\circ} = \operatorname{Im} (u - i). \end{array}
\]

By prop. 7, c) of IV, p. 236, it suffices to show that the image of \(u + i\) (resp. of \(u - i\) ) is closed in E. Since \(u\) is symmetric, we have \(\langle (u + i)(x) \mid (u + i)(y) \rangle = (x \mid y)_u\) for all \(x\) and \(y\) in \(\text{dom}(u)\) . The map \(x \mapsto u(x) + ix\) from \(E_u\) into E is therefore isometric. Its image is closed in E (lemma 8 of I, p. 107). Since this is also the image of \(u + i\) , the latter is closed in E. Likewise, the image of \(u - i\) is closed in E.

Lemma 7. --- Let u be a closed symmetric operator on a complex Hilbert space E and v a closed symmetric extension of u. The domain of u is a closed subspace of the Hilbert space \(E_{v}\) .

Since v is an extension of u, we have \((x \mid y)_{v} = (x \mid y)_{u}\) for x and y in \(\text{dom}(u)\) . The canonical injection of \(E_{u}\) into \(E_{v}\) is therefore isometric, and the conclusion follows (lemma 8 of I, p. 107).

PROPOSITION 26. --- Let u be a closed symmetric operator on a complex Hilbert space E. Let \((\mathrm{E}_{+}, \mathrm{E}_{-})\) be the deficiency pair of u. The spaces \(E_{+}\) , \(E_{-}\) and dom(u) are closed and mutually orthogonal subspaces of \(E_{u^{*}}\) whose Hilbertian sum is equal to \(E_{u^{*}}\) .

Since u is symmetric, we have \(u \subset u^{*}\) , hence the space \(\text{dom}(u)\) is closed in \(E_{u^{*}}\) (lemma 7). The subspaces \(E_{+}\) and \(E_{-}\) are closed in E and contained in \(\text{dom}(u^{*})\) , hence closed in \(E_{u^{*}}\) .

Let \(x \in \mathrm{E}_{+}\) . We have \(u^{*}(x) = ix\) , whence \(\langle x | u(y) \rangle = -i\langle x | y \rangle\) for all \(y \in \operatorname{dom}(u)\) . Consequently, for every \(y \in \operatorname{dom}(u)\) , we have

\[
(x \mid y) _ {u ^ {*}} = \langle x \mid y \rangle + \langle u ^ {*} (x) \mid u ^ {*} (y) \rangle = i (\langle x \mid u (y) \rangle - \langle x \mid u ^ {*} (y) \rangle) = 0
\]

since u is symmetric. The spaces \(E_{+}\) and \(\operatorname{dom}(u)\) are therefore orthogonal in \(E_{u^{*}}\) . Likewise, \(E_{-}\) and \(\operatorname{dom}(u)\) are orthogonal in \(E_{u^{*}}\) .

Let \(x \in \mathrm{E}_{+}\) and \(y \in \mathrm{E}_{-}\) . Then

\[
(x \mid y) _ {u ^ {*}} = \langle x \mid y \rangle + \langle u ^ {*} (x) \mid u ^ {*} (y) \rangle = \langle x \mid y \rangle - \langle x \mid y \rangle = 0,
\]

hence \(\mathrm{E}_{+}\) and \(\mathrm{E}_{-}\) are orthogonal in \(\mathrm{E}_{u^{*}}\) .

Let \(x \in \mathrm{E}_{u^*}\) orthogonal to the closed subspace \(\operatorname{dom}(u) \oplus \mathrm{E}_+ \oplus \mathrm{E}_-\) . For every \(y \in \operatorname{dom}(u)\) , we have \(u^*(y) = u(y)\) since \(u \subset u^*\) , whence

\[
0 = (x \mid y) _ {u ^ {*}} = \langle x \mid y \rangle + \langle u ^ {*} (x) \mid u ^ {*} (y) \rangle = \langle x \mid y \rangle + \langle u ^ {*} (x) \mid u (y) \rangle .
\]

This implies that \(z = u^{*}(x)\) belongs to the domain of \(u^{*}\) and satisfies \(u^{*}(z) = -x\) . Let \(x_{-} = z - ix\) . It is an element of \(\text{dom}(u^{*})\) satisfying \(u^{*}(x_{-}) = -x - iz = -ix_{-}\) . Hence \(x_{-}\) belongs to \(E_{-}\) . But, for every \(w \in E_{-}\) , we have

\[
\begin{array}{c} - i \langle x _ {-} | w \rangle = - i \langle z | w \rangle + \langle x | w \rangle \\ = \langle u ^ {*} (x) | u ^ {*} (w) \rangle + \langle x | w \rangle = (x | w) _ {u ^ {*}} = 0. \end{array}
\]

Taking \(w = x_{-}\) , we deduce that \(x_{-} = 0\) , that is, \(z = ix\) ; since \(z = u^{*}(x)\) , we therefore have \(x \in \mathrm{E}_{+}\) , whence \(x = 0\) since \(x\) is orthogonal to \(\mathrm{E}_{+}\) in \(\mathrm{E}_{u^{*}}\) . This completes the proof that the direct sum \(\operatorname{dom}(u) \oplus \mathrm{E}_{+} \oplus \mathrm{E}_{-}\) is equal to \(\mathrm{E}_{u^{*}}\) .

COROLLARY. --- Let u be a closed symmetric operator on E. Then u is self-adjoint if and only if the deficiency index of u is equal to (0,0).

Indeed, the symmetric operator \(u\) is self-adjoint if and only if \(\text{dom}(u^*) = \text{dom}(u)\) ; but the proposition shows that \(\text{dom}(u^*)\) is the Hilbert sum of \(\text{dom}(u)\) and of the deficiency subspaces of \(u\) .

\pandocbounded{\includegraphics[keepaspectratio]{images/part002_f39b57bf.jpg}}

Let \(u\) be a closed symmetric partial operator. The spectrum of \(u\) is contained in \(\mathbf{R}\) if and only if \(u\) is self-adjoint. Indeed, we know that \(\mathrm{Sp}(u) \subset \mathbf{R}\) if \(u\) is self-adjoint (prop. 17 of IV, p. 248); if \(u\) is not self-adjoint, one of the subspaces \(\mathrm{Ker}(u^{*} + i)\) or \(\mathrm{Ker}(u^{*} - i)\) is nonzero, hence the image of \(u + i\) or of \(u - i\) is a proper subspace of E (prop. 25 of IV, p. 256), so that \(i\) or \(-i\) belongs to \(\mathrm{Sp}(u)\) .

DEFINITION 11. --- Let u be a closed symmetric operator on a complex Hilbert space E. A boundary condition for u is a partially isometric linear map (EVT, V, p. 41, def. 3) from \(\operatorname{Ker}(u^{*} - i)\) into \(\operatorname{Ker}(u^{*} + i)\) .

Let \(u\) be a closed symmetric operator on a complex Hilbert space E and \(b\) a boundary condition for \(u\) . Let \(I = \text{Ker}(b)^{\circ}\) be the initial subspace of \(b\) . We denote by \(u_b\) the reduction of \(u^*\) to the subspace of \(\text{dom}(u^*)\) that is the direct sum of \(\text{dom}(u)\) and the graph in \(E_+ \oplus E_-\) of the restriction of \(b\) to I. Since \(E_{u^*} = \text{dom}(u) \oplus E_+ \oplus E_-\) (prop. 26), the partial operator \(u_b\) is well defined.

In other words, the domain of \(u_{b}\) is the space of elements \(x \in E\) of the form \(x = x_{0} + y + b(y)\) , where \(x_{0} \in \text{dom}(u)\) and y belongs to the initial subspace of b. We then have \(u_{b}(x_{0} + y + b(y)) = u(x_{0}) + iy - ib(y)\) .

THEOREM 1 (von Neumann). --- Let u be a closed symmetric operator on a complex Hilbert space E and \((\mathrm{E}_{+}, \mathrm{E}_{-})\) its deficiency pair.

\begin{enumerate}
\def\labelenumi{\alph{enumi})}
\item
  For every boundary condition b for u, the partial operator \(u_{b}\) is a closed symmetric extension of u ;
\item
  The map \(b \mapsto u_{b}\) is a bijection from the set of boundary conditions for u onto the set of closed symmetric extensions of u;
\item
  For every boundary condition b for u, we have
\end{enumerate}

\[
\begin{array}{l} \operatorname{Im} (u _ {b} + i) = \operatorname{Im} (u + i) \oplus \operatorname{Im} (b) \\ \operatorname{Im} (u _ {b} - i) = \operatorname{Im} (u - i) \oplus \operatorname{Ker} (b) ^ {\circ}; \end{array}
\]

\begin{enumerate}
\def\labelenumi{\alph{enumi})}
\setcounter{enumi}{3}
\tightlist
\item
  For every boundary condition b for u, the deficiency pair of \(u_{b}\) is \((\operatorname{Ker}(b), \operatorname{Ker}(b^{*}))\) .
\end{enumerate}

Let \(b\) be a boundary condition for \(u\) and \(\mathrm{I} \subset \mathrm{E}_{+}\) its initial subspace. The restriction of \(u^{*}\) to \(\Gamma_{b|I} \subset \mathrm{E}_{+} \oplus \mathrm{E}_{-}\) is the continuous linear map defined by \(x + b(x) \mapsto ix - ib(x)\) for \(x \in I\) . The graph of \(u_{b}\) is the direct sum of the graph of u and the graph of this linear map; it is therefore closed. The closed operator \(u_{b}\) is an extension of u such that \(u_{b} \subset u^{*}\) .

Let \(\gamma_{1}\) and \(\eta_{1}\) be elements of I. Consider the elements

\[
\gamma = \gamma_ {1} + b (\gamma_ {1}), \quad \eta = \eta_ {1} + b (\eta_ {1})
\]

of the graph \(\Gamma_{b|I}\) . We compute

\[
\begin{array}{c} \langle \gamma | u ^ {*} (\eta) \rangle = i \big (\langle b (\gamma_ {1}) | \eta_ {1} \rangle - \langle \gamma_ {1} | b (\eta_ {1}) \rangle \big) + i \big (\langle \gamma_ {1} | \eta_ {1} \rangle - \langle b (\gamma_ {1}) | b (\eta_ {1}) \rangle \big) \\ = i \big (\langle b (\gamma_ {1}) | \eta_ {1} \rangle - \langle \gamma_ {1} | b (\eta_ {1}) \rangle \big) \end{array}
\]

since b is isometric on I. From this we deduce that we have

\[
\langle u ^ {*} (\gamma) \mid \eta \rangle = \langle \gamma \mid u ^ {*} (\eta) \rangle\tag{12}
\]

for all \(\gamma\) and \(\eta\) in \(\Gamma_{b|I}\) .

Let x and y be elements of \(\operatorname{dom}(u)\) , \(\gamma\) and \(\eta\) be elements of \(\Gamma_{b|I}\) . We obtain

\[
\begin{array}{c} \langle x + \gamma | u _ {b} (y + \eta) \rangle = \langle x | u (y) \rangle + \langle x | u ^ {*} (\eta) \rangle + \langle \gamma | u (y) \rangle + \langle \gamma | u ^ {*} (\eta) \rangle \\ = \langle u (x) | y \rangle + \langle u (x) | \eta \rangle + \langle u ^ {*} (\gamma) | y \rangle + \langle u ^ {*} (\gamma) | \eta \rangle \\ = \langle u _ {b} (x + \gamma) | y + \eta \rangle , \end{array}
\]

where we used the fact that u is symmetric and that \(\gamma\) and \(\eta\) belong to the domain of \(u^{*}\) as well as formula (12). The operator \(u_{b}\) is therefore a closed symmetric extension of u. Assertion a) is proved.

The map \(b \mapsto u_b\) is injective. Indeed, a partially isometric map from \(E_+\) into \(E_-\) is uniquely determined by its initial subspace and by its restriction to that subspace. Now the domain of \(u_b\) determines the initial subspace I of \(b\) and the graph of the restriction of \(b\) to I.

Let us prove that the map \(b \mapsto u_b\) is surjective. Let \(w\) be a closed symmetric extension of \(u\) . We have \(w \subset w^* \subset u^*\) . The domain of \(w\) is therefore a subspace of \(\text{dom}(u^*)\) containing \(\text{dom}(u)\) , and \(w\) is the reduction of \(u^*\) to this subspace. Let G be the intersection of the domain of \(w\) and of \(E_+ \oplus E_-\) . This is a closed subspace of \(E_{u^*}\) (Lemma 7) and we have \(\text{dom}(w) = \text{dom}(u) \oplus G\) by proposition 26.

Let \(x \in \mathrm{E}_{+}\) and \(y \in \mathrm{E}_{-}\) such that \(x + y \in \mathrm{G}\) . Since \(w\) is symmetric and is the reduction of \(u^{*}\) to \(\operatorname{dom}(w)\) , we have

\[
\langle x + y \mid u ^ {*} (x + y) \rangle = \langle u ^ {*} (x + y) \mid x + y \rangle .
\]

This equality is equivalent to

\[
\begin{array}{c} \langle x | i x \rangle + \langle x | - i y \rangle + \langle y | i x \rangle + \langle y | - i y \rangle \\ = \langle i x | x \rangle + \langle i x | y \rangle + \langle - i y | x \rangle + \langle - i y | y \rangle , \end{array}
\]

that is, to \(\|x\|^{2} = \|y\|^{2}\) . In particular, the canonical projection \(p_{+}: G \to E_{+}\) is injective, and if I denotes its image, the space G is the graph of an isometric map \(b_{0}\) from I into \(E_{-}\) . For x in I and y in \(E_{-}\) such that \(x + y \in G\) , we have

\[
\begin{array}{r} \| x + y \| _ {u ^ {*}} ^ {2} = \| x \| _ {u ^ {*}} ^ {2} + \| y \| _ {u ^ {*}} ^ {2} = 2 (\| x \| ^ {2} + \| y \| ^ {2}) \\ = 4 \| x \| ^ {2} = 4 \| p _ {+} (x + y) \| ^ {2}. \end{array}
\]

Thus, considering G as a closed subspace of the Hilbert space \(E_{u^{*}}\) , the map \(x \mapsto \frac{1}{2}p_{+}(x)\) is an isometry from G onto I. In particular, I is closed in \(E_{+}\) (lemma 8 of I, p. 107). There then exists a unique boundary condition b for u for which I is the initial space and which coincides with \(b_{0}\) on I. We have \(\text{dom}(w) = \text{dom}(u_{b})\) , whence \(w = u_{b}\) . This establishes assertion b).

Let us prove the assertions \(c\) ) and \(d\) ). Let \(b\) be a boundary condition for \(u\) and I its initial subspace. For \(x_0 \in \mathrm{dom}(u)\) and \(y \in \mathrm{I}\) , we have

\[
u _ {b} (x _ {0} + y + b (y)) - i (x _ {0} + y + b (y)) = (u - i) (x _ {0}) - 2 i b (y).
\]

One has \(\operatorname{Im}(u - i) = \operatorname{E}_{-}^{\circ}\) (prop. 25). Since the image of \(b\) is contained in \(\operatorname{E}_{-}\) , this formula shows that \(\operatorname{Im}(u_b - i) = \operatorname{Im}(u - i) \oplus \operatorname{Im}(b)\) . By loc. cit., we therefore have \(\operatorname{Ker}(u_b^* + i) = \operatorname{E}_{-} \cap \operatorname{Im}(b)^{\circ} = \operatorname{Ker}(b^*)\) .

Similarly, one verifies that \(\operatorname{Im}(u_b + i) = \operatorname{Im}(u + i) \oplus \operatorname{I}\) and hence (loc. cit.) that \(\operatorname{Ker}(u_b^* - i) = \operatorname{I}^\circ \cap \operatorname{E}_+ = \operatorname{Ker}(b)\) .

COROLLARY 1. --- Let u be a closed symmetric operator on a complex Hilbert space E. The map \(b \mapsto u_b\) induces a bijection between the set of isometric isomorphisms from \(\text{Ker}(u^* - i)\) onto \(\text{Ker}(u^* + i)\) and the set of self-adjoint extensions of u. In particular, there exists a self-adjoint extension of u if and only if the two components of the deficiency index of u are equal.

This follows from the corollary of Proposition 26 and the theorem, as well as from EVT, V, p. 25, cor. 2.

COROLLARY 2. --- Let E be a real Hilbert space. Let u be a closed symmetric operator on E. The partial operator \(u_{(\mathbf{C})}\) on \(\mathrm{E}_{(\mathbf{C})}\) is symmetric and closed, and admits at least one self-adjoint extension.

The partial operator \(u_{(\mathbf{C})}\) is closed and symmetric by prop. 2 of IV, p. 228 and lemma 3 of IV, p. 239. The R-linear isomorphism \(j\) from \(\mathrm{E}_{(\mathbf{C})}\) to \(\mathrm{E}_{(\mathbf{C})}\) such that \(j(z \otimes x) = \overline{z} \otimes x\) for all \(z \in \mathbf{C}\) and every \(x \in \mathrm{E}\) induces an isomorphism of Hilbert spaces from \(\mathrm{Ker}(u_{(\mathbf{C})}^* - i)\) onto \(\mathrm{Ker}(u_{(\mathbf{C})}^* + i)\), and we can apply Corollary 1.

COROLLARY 3. --- Let u be a symmetric operator on E. The following conditions are equivalent:

\begin{enumerate}
\def\labelenumi{(\roman{enumi})}
\item
  The operator u is essentially self-adjoint;
\item
  The spaces \(\operatorname{Ker}(u^{*} + i)\) and \(\operatorname{Ker}(u^{*} - i)\) are zero;
\item
  The spaces \(\operatorname{Im}(u + i)\) and \(\operatorname{Im}(u - i)\) are dense in E;
\item
  The partial symmetric operator u has a unique self-adjoint extension.
\end{enumerate}

Conditions (ii) and (iii) are equivalent by Proposition 7, c) of IV, p. 236.

The partial operator \(u\) is essentially self-adjoint if and only if \(\overline{u}\) is self-adjoint, that is, if \(\text{Ker}(\overline{u}^* - i) = \text{Ker}(\overline{u}^* + i) = \{0\}\) (corollary of prop. 26). Since \(\overline{u}^* = u^*\) (prop. 8 of IV, p. 237), condition (i) is therefore equivalent to condition (ii). Moreover, the preceding corollary shows that (ii) implies that \(u\) has a unique self-adjoint extension, which is condition (iv).

Finally, suppose that \(u\) has a unique self-adjoint extension \(v\) . Then the same is true of \(\overline{u}\) , since every self-adjoint extension of \(\overline{u}\) is one of \(u\) , hence is equal to \(v\) . By the preceding corollary, the spaces \(\mathrm{Ker}(\overline{u}^{*} + i)\) and \(\mathrm{Ker}(\overline{u}^{*} - i)\) must be zero, whence condition (ii) using again \(\overline{u}^{*} = u^{*}\) .

Example. --- Let U be an open subset of \(\mathbf{R}^{n}\) endowed with Lebesgue measure \(\mu\) . Let \(\Delta\) be the scalar differential operator

\[
\Delta = - \sum_ {i = 1} ^ {n} \partial_ {i} ^ {2}
\]

on U. Denote by \(u\) the partial operator on the real Hilbert space \(\mathrm{L}_{\mathbf{R}}^{2}(\mathrm{U},\mu)\) with domain \(\mathcal{D}_{\mathbf{R}}(\mathrm{U})\) (IV, p. 201, remark) defined by \(\varphi \mapsto -\sum_{i=1}^{n}\partial_{i}^{2}\varphi\) . We have \(u_{(\mathbf{C})} = \Delta_{-}\) , which is a closable partial operator (prop. 13 of IV, p. 242) and symmetric (since \(\Delta\) is formally symmetric), hence \(u\) is closable and symmetric (proposition 2 of IV, p. 228 and lemma 3 of IV, p. 239). By the above corollary, there therefore exists a self-adjoint extension of \(u_{(\mathbf{C})}\) . This is a Laplacian on U (IV, p. 243).

\section{NORMAL PARTIAL OPERATORS AND THE SPECTRAL THEOREM}

\subsection{Boundedification}

Let E be a complex Hilbert space. We denote by s the unitary map \((x,y)\mapsto(-y,x)\) on \(E\oplus E\) . Let u be a closed partial operator with dense domain on E. The partial operator \(u^{*}u\) is self-adjoint and positive (prop. 12 of IV, p. 241), hence \(-1\notin\mathrm{Sp}(u^{*}u)\) (prop. 17 of IV, p. 248). We denote by \(\mathrm{W}(u)=(1_{\mathrm{E}}+u^{*}u)^{-1}=-\mathrm{R}(u^{*}u,-1)\) ; this is a positive injective endomorphism of E.

Denote by \(p_{1}\) and \(p_{2}\) the two canonical projections of \(\Gamma_{u}\) into E. These are elements of \(\mathcal{L}(\Gamma_{u};\mathrm{E})\) , and we have the equality of correspondences \(p_{2}=u\circ p_{1}\) . Let \((j,|p_{1}^{*}|)\) be the polar decomposition (def. 4 of I, p. 140) of the endomorphism \(p_{1}^{*}\in\mathcal{L}(\mathrm{E};\Gamma_{u})\) , so that \(p_{1}^{*}=|p_{1}^{*}|\circ j\) . The map j is a partial isometry from E into \(\Gamma_{u}\) .

DEFINITION 1. --- The endomorphism \(p_2 \circ j\) of E is called the boundedification of \(u\) ; we denote by \(b(u)\) this endomorphism.

One has \(\|b(u)\|\leqslant1\) .

PROPOSITION 1. --- We have the formulas

\[
| p _ {1} ^ {*} | = \mathrm{W} (u) ^ {1 / 2},\tag{1}
\]

\[
b (u) = u \circ \mathrm{W} (u) ^ {1 / 2} = u \circ | p _ {1} ^ {*} |,\tag{2}
\]

\[
1 _ {\mathrm{E}} - b (u) ^ {*} b (u) = \mathrm{W} (u),\tag{3}
\]

\[
b (u) \mathrm{W} (u) = \mathrm{W} (u ^ {*}) b (u), \quad b (u) \mathrm{W} (u) ^ {1 / 2} = \mathrm{W} (u ^ {*}) ^ {1 / 2} b (u),\tag{4}
\]

Let \(x \in E\) . Set \(y = \mathrm{W}(u)(x) \in \mathrm{dom}(u^{*}u)\) . We have \(y \in \mathrm{dom}(u)\) and \(u(y) \in \mathrm{dom}(u^{*})\) . For every element \((y_{1}, u(y_{1}))\) of \(\Gamma_{u}\) , where \(y_{1} \in \mathrm{dom}(u)\) , one computes

\[
\begin{array}{c} \langle (y, u (y)) | (y _ {1}, u (y _ {1})) \rangle = \langle y | y _ {1} \rangle + \langle u (y) | u (y _ {1}) \rangle \\ = \langle y + u ^ {*} u (y) | y _ {1} \rangle = \langle x | p _ {1} (y _ {1}, u (y _ {1})) \rangle . \end{array}
\]

This means that \(p_{1}^{*}(x) = (y, u(y))\) , whence \(p_{1} \circ p_{1}^{*}(x) = y = \mathrm{W}(u)(x)\) . Consequently, one has \(p_{1} \circ p_{1}^{*} = \mathrm{W}(u)\) , hence \(|p_{1}^{*}| = \mathrm{W}(u)^{1/2}\) , that is, formula (1). In particular, since \(\operatorname{Im}(p_{1}) = \operatorname{Im}(|p_{1}^{*}|)\) by the cor. of prop. 11 of I, p. 140, it follows that \(\operatorname{dom}(u) = \operatorname{Im}(p_{1}) = \operatorname{Im}(\mathrm{W}(u)^{1/2})\) .

The partial operator \(u \circ \mathrm{W}(u)^{1/2}\) therefore has domain E. The relation \(p_1 \circ j = (p_1^*)^* \circ j = |p_1^*| = \mathrm{W}(u)^{1/2}\) (prop. 11, a) of I, p. 140) then implies that

\[
b (u) = p _ {2} \circ j = u \circ p _ {1} \circ j = u \circ \mathrm{W} (u) ^ {1 / 2},
\]

which is formula (2).

One has \(\operatorname{Ker}(j) = \operatorname{Ker}(|p_1^*|)\) (prop. 10, \(b\) ) of I, p. 139) and formula (1) therefore imply that the partial isometry \(j: \mathrm{E} \to \Gamma_u\) is injective, hence isometric, whence \(j^* \circ j = 1_{\mathrm{E}}\) . Writing \(j = (p_1 \circ j, b(u))\) , we find

\[
1 _ {\mathrm{E}} = j ^ {*} \circ j = (p _ {1} \circ j) ^ {*} (p _ {1} \circ j) + b (u) ^ {*} b (u) = \mathrm{W} (u) + b (u) ^ {*} b (u),
\]

whence formula (3).

Let \(x \in \mathrm{dom}(u)\) and set \(y = \mathrm{W}(u)(x)\) . We have \(y \in \mathrm{dom}(u)\) and the formula \(y + u^{*}u(y) = x\) implies that \(u^{*}u(y) \in \mathrm{dom}(u)\) . We then have

\[
\begin{array}{r l} & {\mathrm{W} (u ^ {*}) ^ {- 1} (u (y)) = (1 _ {\mathrm{E}} + u u ^ {*}) (u (y))} \\ & {\qquad = u (y) + u u ^ {*} u (y) = u (y + u ^ {*} u (y)) = u (x),} \end{array}
\]

which means that the reduction of \(u \circ \mathrm{W}(u)\) to the domain of u is equal to \(\mathrm{W}(u^{*}) \circ u\) . Since the image of \(|p_{1}^{*}|\) is equal to that of \(p_{1}\), and hence to the domain of u (cor. of prop. 11 of I, p. 140), it follows

\[
u \circ \mathrm{W} (u) \circ | p _ {1} ^ {*} | = \mathrm{W} (u ^ {*}) \circ u \circ | p _ {1} ^ {*} | = \mathrm{W} (u ^ {*}) \circ b (u)
\]

(formula (2)). Moreover, \(|p_{1}^{*}| = \mathrm{W}(u)^{1/2}\) commutes with \(\mathrm{W}(u)\) , hence we obtain

\[
b (u) \circ \mathrm{W} (u) = u \circ | p _ {1} ^ {*} | \circ \mathrm{W} (u) = u \circ \mathrm{W} (u) \circ | p _ {1} ^ {*} | = \mathrm{W} (u ^ {*}) \circ b (u).
\]

This relation implies that \(b(u) \circ f(\mathrm{W}(u)) = f(\mathrm{W}(u^{*})) \circ b(u)\) for every function \(f \in \mathcal{C}(\mathbf{R}_{+})\) (prop. 11 of I, p. 113), hence in particular that \(b(u) \circ \mathrm{W}(u)^{1/2} = \mathrm{W}(u^{*})^{1/2} \circ b(u)\) . This proves formula (4) and concludes the proof.

COROLLARY. --- We have \(b(u)^{*} = b(u^{*})\) .

Denote by \(q_{1}\) and \(q_{2}\) the projections \(\Gamma_{u^{*}}\to\mathrm{E}\) and let \((k,|q_{1}^{*}|)\) be the polar decomposition of \(q_{1}^{*}\) , so that \(b(u^{*})=q_{2}\circ k\) . For all \(x\) and \(y\) in E, we have \(j(x)\in\Gamma_{u}\) and \(k(y)\in\Gamma_{u^{*}}\) , and since \(\Gamma_{u}\) is orthogonal to \(s(\Gamma_{u^{*}})\) by prop. 7 of IV, p. 236, it follows

\[
0 = \langle j (x) | s (k (y)) \rangle = \langle p _ {1} \circ j (x) | - q _ {2} \circ k (y) \rangle + \langle p _ {2} \circ j (x) | q _ {1} \circ k (y) \rangle .
\]

Since \(p_1 \circ j = |p_1^*|\) and \(q_1 \circ k = |q_1^*|\) (prop. 11, b) of I, p. 140) we obtain

\[
\langle b (u) x \mid | q _ {1} ^ {*} | y \rangle = \langle | p _ {1} ^ {*} | x \mid b (u ^ {*}) y \rangle ,
\]

hence \(|q_1^*|b(u) = b(u^*)^* |p_1^*|\) . Using formulas (1) and (4), we conclude from this that \(b(u)|p_1^*| = |q_1^*|b(u) = b(u^*)^* |p_1^*|\) , and since the image of \(|p_1^*|\) is the domain of \(u\) which is dense in E, we deduce \(b(u) = b(u^*)^*\) .

We denote \(\Omega(\mathrm{E})\) the set of \(v \in \mathcal{L}(\mathrm{E})\) such that \(1_{\mathrm{E}} - v^{*}v\) is positive and injective. For \(v \in \Omega(\mathrm{E})\) , the Hermitian endomorphism \((1_{\mathrm{E}} - v^{*}v)^{1/2}\) is injective, since its square is, and we denote by \(\mathrm{B}(v)\) the partial operator \(v \circ ((1_{\mathrm{E}} - v^{*}v)^{1/2})^{-1}\) .

Note that \(1_{\mathrm{E}} - v^{*}v\) is positive if and only if \(\| v\| \leqslant 1\) , since we have \(\langle x|(1_{\mathrm{E}} - v^{*}v)(x)\rangle = \| x\| ^2 -\| v(x)\| ^2\) for all \(x\in \mathrm{E}\) .

Lemma 1. --- The subset \(\Omega(\mathrm{E})\) is self-adjoint in \(\mathcal{L}(\mathrm{E})\) .

Let \(v \in \Omega(\mathrm{E})\) . We have \(\|v^{*}\| = \|v\| \leqslant 1\) . Therefore the endomorphism \(1_{\mathrm{E}} - vv^{*}\) is positive. It is injective: if \(x \in \operatorname{Ker}(1_{\mathrm{E}} - vv^{*})\) , we have \(vv^{*}(x) = x\) , whence \(v^{*}(v(v^{*}(x))) = v^{*}(x)\) , then \(v^{*}(x) = 0\) since \(1_{\mathrm{E}} - v^{*}v\) is injective, and finally \(x = v(v^{*}(x)) = 0\) . The lemma follows from this.

PROPOSITION 2. — The map \(u \mapsto b(u)\) defines a bijection from the set of closed densely defined partial operators on E onto the set \(\Omega(E)\) . The inverse bijection is given by \(v \mapsto B(v)\) .

The relation \(1_{\mathrm{E}} - b(u)^{*}b(u) = \mathrm{W}(u)\) (formula (3)) implies that \(b(u)\) belongs to \(\Omega (\mathrm{E})\) since \(\mathrm{W}(u)\) is positive and injective. Moreover, since \(b(u) = u\circ \mathrm{W}(u)^{1 / 2}\) (formula (2)), we obtain

\[
u = b (u) \circ (\mathrm{W} (u) ^ {1 / 2}) ^ {- 1} = b (u) \circ ((1 - b (u) ^ {*} b (u)) ^ {1 / 2}) ^ {- 1} = \mathrm{B} (b (u)).
\]

Conversely, let \(v \in \Omega(\mathrm{E})\) . Set \(w = (1_{\mathrm{E}} - v^{*}v)^{1/2}\) and set \(u = \mathrm{B}(v) = v \circ w^{-1}\) . The domain of u is the image of w, which is dense in E since w is Hermitian and injective (EVT, V, p. 41, prop. 4 (i)). For every \(x \in \mathrm{dom}(u)\) , it follows that

\[
\begin{array}{r l} & {\| u (x) \| ^ {2} = \langle (v ^ {*} v) (w ^ {- 1} (x)) | w ^ {- 1} (x) \rangle} \\ & {\qquad = - \langle (1 _ {\mathrm{E}} - v ^ {*} v) (w ^ {- 1} (x)) | w ^ {- 1} (x) \rangle + \| w ^ {- 1} (x) \| ^ {2}} \\ & {\qquad = - \langle w (x) | w ^ {- 1} (x) \rangle + \| w ^ {- 1} (x) \| ^ {2}} \end{array}
\]

since \((1_{\mathrm{E}} - v^{*}v) \circ w^{-1}\) is the reduction of w to the domain of \(w^{-1}\) . Since w is self-adjoint, it follows that

\[
\| w ^ {- 1} (x) \| ^ {2} = \| x \| ^ {2} + \| v \circ w ^ {- 1} (x) \| ^ {2} = \| x \| ^ {2} + \| u (x) \| ^ {2},
\]

and the partial operator \(\mathrm{B}(v) = u = v \circ w^{-1}\) is therefore closed by Lemma 2 of IV, p. 231.

Let us finally prove that \(b(\mathrm{B}(v)) = v\) . Since \(u = v \circ w^{-1}\) , we have \(v = u \circ w\) and it therefore suffices to check that \(w = \mathrm{W}(u)^{1/2}\) by formula (2), or even that \(1_{\mathrm{E}} - v^{*}v = \mathrm{W}(u)\) .

The endomorphism \(v^*\) belongs to \(\Omega(\mathrm{E})\) . By lemma 1, set \(w' = (1_{\mathrm{E}} - vv^{*})^{1/2}\) . It is an injective positive endomorphism of E, and we have \(v \circ w = w' \circ v\) (prop. 11 of I, p. 113). The graph of \(u\) is the set of elements of the form \((w(x), v(x))\) for \(x \in \mathrm{E}\) . By prop. 7 of IV, p. 236, the graph of \(u^*\) is \(s(\Gamma_u)^{\circ}\) . Since, for all \(x\) and \(y \in \mathrm{E}\) , we have

\[
\langle (w ^ {\prime} (x), v ^ {*} (x)) | (- v (y), w (y)) \rangle = - \langle x | w ^ {\prime} v (y) \rangle + \langle x | v w (y) \rangle = 0,
\]

the graph of \(u^*\) contains the elements of the form \((w'(x), v^*(x))\) for \(x \in \mathrm{E}\) , and we then have \(u^*(w'(x)) = v^*(x)\) . In particular, the domain of \(u^*\) contains the image of \(w'\) .

Let \(x \in E\) and set \(y = w^{2}(x) = (1_{\mathrm{E}} - v^{*}v)(x)\) , hence \(x = y + v^{*}v(x)\) . We have \(y \in \operatorname{dom}(u)\) and \(u(y) = v(w^{-1}(y)) = v(w(x)) = w'(v(x))\) . In particular, \(u(y) \in \operatorname{Im}(w) \subset \operatorname{dom}(u^{*})\) , and \(u^{*}(u(y)) = v^{*}(v(x))\) . It follows therefore \(x = y + v^{*}v(x) = y + u^{*}u(y)\) , whence \(y = \operatorname{W}(u)(x)\) , that is,

\[
(1 _ {\mathrm{E}} - v ^ {*} v) (x) = \mathrm{W} (u) (x).
\]

We have therefore shown that \(1_{\mathrm{E}} - v^{*}v = \mathrm{W}(u)\) , as desired.

COROLLARY. --- The endomorphism \(b(u)\) is Hermitian if and only if u is self-adjoint.

This follows from the injectivity of the map \(u \mapsto b(u)\) (prop. 2) and from the formula \(b(u^{*}) = b(u)^{*}\) (cor. of prop. 1).

\subsection{Normal partial operators and the spectral theorem}

In this section, E is a complex Hilbert space.

DEFINITION 2. --- Let u be a partial operator on E. One says that u is normal if u is closed with dense domain and if its bounded transform b(u) is a normal endomorphism of E.

If \(u \in \mathcal{L}(\mathrm{E})\) this definition coincides with EVT, V, p. 42, def. 4, by virtue of the formulas \(1_{\mathrm{E}} - b(u)^{*}b(u) = \mathrm{W}(u) = (1_{\mathrm{E}} + u^{*}u)^{-1}\) (prop. 1 of IV, p. 262) and \(b(u^{*}) = b(u)^{*}\) (cor. of loc. cit.).

If \(u\) is a self-adjoint partial operator on \(\mathrm{E}\) , then \(b(u)\) is Hermitian (cor. of prop. 2 of IV, p. 264) hence \(u\) is normal.

Let D be the open unit disk in C. We denote by \(\beta\) the function from C to D defined by \(\beta(z)=z/\sqrt{1+|z|^{2}}\) . It is a homeomorphism, whose inverse satisfies \(\beta^{-1}(z)=z/\sqrt{1-|z|^{2}}\) for \(z\in D\) .

Let \(u \in \mathcal{L}(\mathrm{E})\) . It follows from formulas (2) and (3) of proposition 1 of IV, p. 262 that \(u = \beta^{-1}(b(u))\) , and hence

\[
b (u) = \beta (u).\tag{5}
\]

Lemma 2. --- Let X be a locally compact topological space and let μ be a positive measure on X.

\begin{enumerate}
\def\labelenumi{\alph{enumi})}
\item
  Let \(g\) be a function \(\mu\)-measurable on X. The multiplication operator \(m_g\) on \(\mathrm{L}^2 (\mathrm{X},\mu)\) is normal and \(b(m_g) = m_{\beta \circ g}\) ;
\item
  Let \(h\colon X\to D\) be a function \(\mu\)-measurable. The endomorphism \(m_h\) belongs to \(\Omega(L^2(X,\mu))\) and \(B(m_h)=m_{\beta^{-1}\circ h}\) .
\end{enumerate}

The partial operator \(m_g\) is closed with dense domain. Since \(m_g^* m_g = m_{\overline{g}} m_g = m_{|g|^2}\) (prop. 23 of IV, p. 253 and prop. 24 of IV, p. 255) we have

\[
\mathrm{W} (m _ {g}) = - \mathrm{R} (m _ {| g | ^ {2}}, - 1) = m _ {(1 + | g | ^ {2}) ^ {- 1}}
\]

(prop. 22 of IV, p. 252). This implies

\[
b (m _ {g}) = m _ {g} \circ \mathrm{W} (m _ {g}) ^ {1 / 2} = m _ {g} \circ m _ {(1 + | g | ^ {2}) ^ {- 1 / 2}} = m _ {\beta \circ g}
\]

according to formula (2) of IV, p. 262, the cor. of prop. 6 of IV, p. 187 and lemma 6 of IV, p. 254. Finally applying prop. 6 of IV, p. 187, we obtain assertion a).

Let us prove \(b\) ). Since \(1 - |h(x)|^2 > 0\) for all \(x \in \mathrm{X}\) , we have \(m_h \in \Omega(\mathrm{L}^2(\mathrm{X}, \mu))\) (cor. of prop. 6 of IV, p. 187). Since the map \(u \mapsto b(u)\) is injective (prop. 2 of IV, p. 264) and the endomorphisms \(b(\mathrm{B}(m_h))\) and \(b(m_{\beta^{-1}\circ h})\) are equal to \(m_h\) according to assertion \(a\) ), we indeed have \(\mathrm{B}(m_h) = m_{\beta^{-1}\circ h}\) .

THEOREM 1 ("Spectral theorem"). --- Let u be a normal partial operator on the Hilbert space E. There exist a locally compact topological space X, a positive measure μ on X, an isometric isomorphism θ of L²(X,μ) onto E and a continuous function g on X, such that \(u = \theta \circ m_g \circ \theta^{-1}\).

The endomorphism \(b(u)\) of E is normal by definition. By the spectral theorem for normal endomorphisms of E (corollary 1 of IV, p. 193), there therefore exist a locally compact topological space \(\widetilde{X}\), a positive measure \(\mu\) on \(\widetilde{X}\), an isometric isomorphism \(\widetilde{\theta}\) of \(\mathrm{L}^2 (\widetilde{\mathrm{X}},\mu)\) onto E, and a bounded continuous function \(h\) on \(\widetilde{X}\), so that \(b(u)\) coincides with \(\widetilde{\theta}\circ m_h\circ \widetilde{\theta}^{-1}\) .

Let N be the closed subset of \(x \in \widetilde{X}\) such that \(|h(x)| \geqslant 1\) . Since the endomorphism \(1 - b(u)^{*}b(u) = \widetilde{\theta} \circ m_{1-|h|^2} \circ \widetilde{\theta}^{-1}\) is positive and injective, the set N is locally \(\mu\) -negligible (lemma 7 of IV, p. 186 and cor. of prop. 6 of IV, p. 187). We then set X = \(\widetilde{X} - N\) . This is a locally compact space, and the restriction of functions to X induces an isometric isomorphism of \(\mathrm{L}^{2}(\widetilde{\mathrm{X}}, \mu)\) onto \(\mathrm{L}^{2}(\mathrm{X}, \mu|\mathrm{X})\) (prop. 1 of IV, p. 179). The composition of \(\widetilde{\theta}\) and of the restriction of functions to X therefore induces an isometric isomorphism \(\theta\) of \(\mathrm{L}^{2}(\mathrm{X}, \mu|\mathrm{X})\) onto E. Since \(b(u) = \theta \circ m_{h|\mathrm{X}} \circ \theta^{-1}\) , it follows that \(u = \theta \circ m_{\beta^{-1} \circ (h|\mathrm{X})} \circ \theta^{-1}\) (lemma 2). The theorem follows from this formula.

Lemma 3. --- Let u be a normal partial operator on E. We have

\[
\operatorname{Sp} (b (u)) \cap \mathrm{D} = \beta (\operatorname{Sp} (u)).
\]

One has \(b(u) = u \circ \mathrm{W}(u)^{1/2}\) (prop. 1 of IV, p. 262). The endomorphism \(\mathrm{W}(u)^{1/2}\) is injective and its image is the domain of u (formula (1) of IV, p. 262).

Let \(\lambda \in \mathbf{C}\) . The number \(\lambda\) belongs to the resolvent set of \(u\) if and only if \((u - \lambda 1_{\mathrm{E}}) \circ \mathrm{W}(u)^{1/2}\) is a bijective linear map from E onto E (remark 3 of IV, p. 245). Now, by formula (3) of IV, p. 262, we have

\[
(u - \lambda 1 _ {\mathrm{E}}) \circ \mathrm{W} (u) ^ {1 / 2} = b (u) - \lambda (1 _ {\mathrm{E}} - b (u) ^ {*} b (u)) ^ {1 / 2} = f _ {\lambda} (b (u))
\]

where \(f_{\lambda}\) is the continuous function defined on \(\overline{\mathrm{D}}\) by \(z \mapsto z - \lambda (1 - |z|^2)^{1/2}\) . The endomorphism \(f_{\lambda}(b(u))\) is bijective if and only if its spectrum does not contain 0. Since \(\mathrm{Sp}(f_{\lambda}(b(u))) = f_{\lambda}(\mathrm{Sp}(b(u)))\) (cor. 2 of prop. 7 of I, p. 111), this is the case if and only if 0 does not belong to the set \(f_{\lambda}(\mathrm{Sp}(b(u)))\) . Thus, \(\lambda \in \mathrm{Sp}(u)\) if and only if there exists \(z \in \mathrm{Sp}(b(u)) \cap \overline{\mathrm{D}}\) such that

\[
z - \lambda (1 - | z | ^ {2}) ^ {1 / 2} = 0.
\]

This equality, if valid, implies that \(z \in D\) , and means that \(\lambda = \beta^{-1}(z)\) . We conclude that \(\operatorname{Sp}(u) = \beta^{-1}(\operatorname{Sp}(b(u)) \cap D)\) , as stated.

Let u be a partial normal operator on E. Let \(f \in \mathcal{K}(\mathrm{Sp}(u))\) . The map \(z \mapsto f(\beta^{-1}(z))\) from \(\mathrm{Sp}(b(u)) \cap D\) into C is continuous and compactly supported; the unique mapping \(f_{\beta}\) from \(\mathrm{Sp}(b(u))\) into C which extends it by zero is therefore continuous.

DEFINITION 3. --- For every function \(f \in \mathcal{K}(\mathrm{Sp}(u))\) , we define the endomorphism \(f(u)\) of E by \(f(u) = f_{\beta}(b(u))\) .

The map \(f \mapsto f_{\beta}\) is a morphism of complex algebras from \(\mathcal{K}(\mathrm{Sp}(u))\) into \(\mathcal{C}(\mathrm{Sp}(b(u)))\) , so the map \(f \mapsto f(u)\) is a morphism of complex algebras from \(\mathcal{K}(\mathrm{Sp}(u))\) into \(\mathcal{L}(\mathrm{E})\) . We have \(\overline{f}(u) = f(u)^*\) since \(\overline{f}_{\beta} = \overline{f_{\beta}}\) . If \(f \geqslant 0\) then \(f_{\beta} \geqslant 0\) , hence \(f(u)\) is a positive endomorphism of E.

Remark. --- Suppose that u is a normal endomorphism of E, so that \(\mathrm{Sp}(u)\) is a compact subset of C. Since the bounded transform \(b(u)\) coincides with \(\beta(u)\) (formula (5)), and since \(\mathrm{Sp}(b(u)) = \beta(\mathrm{Sp}(u))\) (cor. 2 of prop. 7 of I, p. 111), the spectrum of \(b(u)\) is a compact subset of D. For every continuous function \(f \in \mathcal{C}(\mathrm{Sp}(u))\) , the function \(f_{\beta}\) coincides with the restriction of \(f \circ \beta^{-1}\) to \(\mathrm{Sp}(u)\) . Consequently, \(f_{\beta}(\beta(u))\) coincides with the endomorphism \(f(u)\) defined by the continuous functional calculus of u. The above definition is therefore compatible with that of the continuous functional calculus of the star algebra \(\mathcal{L}(\mathrm{E})\) .

Let \(f\in \mathcal{K}(\mathrm{Sp}(u))\) . We have

\[
\| f (u) \| \leqslant \| f _ {\beta} \| _ {\infty} = \| f \| _ {\infty},
\]

so that, for all \(x\) and \(y\) in E, we obtain the majorization \(|\langle x | f(u)y\rangle| \leqslant \|x\|\|y\|\|f\|_{\infty}\) . The map \(f \mapsto \langle x | f(u)y\rangle\) is therefore a bounded measure on \(\mathrm{Sp}(u)\) of total mass \(\leqslant \|x\|\|y\|\) (INT, IV, p. 154, § 4, no. 7). If \(x = y\) it is a positive measure, since \(f(u)\) is positive when \(f \geqslant 0\) .

DEFINITION 4. --- Let u be a normal partial operator on a complex Hilbert space E. Let x and y be in E. The bounded measure on Sp(u) defined by \(f \mapsto \langle x | f(u)y \rangle\) for \(f \in \mathcal{K}(\mathrm{Sp}(u))\) is called the spectral measure of \((x,y)\) relative to \(u\) . When \(x = y\) , we say that it is the spectral measure of \(x\) relative to \(u\) .

Remarks. --- 1) When u is continuous, this definition coincides with that of def. 2 of IV, p. 190.

\begin{enumerate}
\def\labelenumi{\arabic{enumi})}
\setcounter{enumi}{1}
\tightlist
\item
  Let \(j\) be the canonical inclusion of \(\mathrm{Sp}(u)\) into \(\mathbf{C}\) . Since \(\mathrm{Sp}(u)\) is closed, the map \(j\) is \(\mu\) -proper for every bounded measure \(\mu\) on \(\mathrm{Sp}(u)\) (INT, V, p. 68, § 6, no. 1, prop. 1). One often identifies the spectral measures of \(u\) with the measures on \(\mathbf{C}\) which are their images under \(j\) (cf. INT, V, p. 84, § 7, no. 2).
\end{enumerate}

The map from E×E into the Banach space \(\mathcal{M}^{1}(\mathrm{Sp}(u))\) that associates to \((x,y)\) the spectral measure of \((x,y)\) relative to u is sesquilinear and continuous with norm \(\leqslant 1\) .

Lemma 4. --- Let X be a locally compact topological space, let \(\mu\) be a positive measure on X and let \(g\) be a function \(\mu\)-measurable on X. Let \(m_g\) be the multiplication operator by \(g\) on \(\mathrm{L}^2 (\mathrm{X},\mu)\) .

\begin{enumerate}
\def\labelenumi{\alph{enumi})}
\item
  The application \(f\mapsto f(m_{g})\) from \(\mathcal{K}(\mathrm{Sp}(m_g))\) into \(\mathcal{L}(\mathrm{L}^{2}(\mathrm{X},\mu))\) is given by \(f\mapsto m_{f\circ g}\) ;
\item
  For \(f_1\) and \(f_2\) in \(\mathcal{L}^2(\mathrm{X},\mu)\) with classes \(\widetilde{f}_1\) and \(\widetilde{f}_2\) in \(\mathrm{L}^2(\mathrm{X},\mu)\) the spectral measure of \((\widetilde{f}_1,\widetilde{f}_2)\) relative to \(m_g\) is the restriction to \(\mathrm{Sp}(m_g)\) of the image measure \(g(\overline{f}_1f_2\cdot \mu)\) .
\end{enumerate}

The partial operator \(m_g\) is normal and \(b(m_g) = m_{\beta \circ g} = m_{g(1+|g|^2)^{-1/2}}\) (lemma 2). Moreover, \(\beta(g(x))\) belongs to \(\mathrm{Sp}(m_{\beta \circ g})\) for all \(x\) outside a locally \(\mu\) -negligible \(Y \subset X\) (proposition 22 of IV, p. 252 and lemma 1 of IV, p. 181). Since \(f_\beta(\beta(g(x))) = f(g(x))\) for all \(x \in X - Y\) , we deduce that \(f(m_g) = m_{f \circ g}\) by definition 3 and the corollary of Proposition 6 of IV, p. 187.

Let \(f_1\) and \(f_2\) be in \(\mathcal{L}^2(\mathrm{X},\mu)\) with classes \(\widetilde{f}_1\) and \(\widetilde{f}_2\) in \(\mathrm{L}^2(\mathrm{X},\mu)\) . Since the measure \(\nu = \overline{f}_1f_2\cdot \mu\) is bounded, the image measure \(g(\nu)\) is defined (INT, V, p. 69, § 6, no. 1, remark 1). Let \(f\in \mathcal{K}(\mathrm{Sp}(m_g))\) . We have \(f(m_{g}) = m_{f\circ g}\) according to \(a\) ), whence

\[
\langle \widetilde {f} _ {1} \mid f (m _ {g}) \widetilde {f} _ {2} \rangle = \int_ {\mathrm{X}} \overline {{f}} _ {1} (f \circ g) f _ {2} d \mu = \int_ {\mathrm{X}} (f \circ g) (\overline {{f}} _ {1} f _ {2} d \mu) = \int_ {\mathbf {C}} f d \nu
\]

(INT, V, p. 69, § 6, n° 1, formula (1)), which establishes assertion b).

Example. --- Let \(n \in \mathbf{N}\) . We equip \(\mathbf{R}^{n}\) with Lebesgue measure, and we identify \(\mathbf{R}^{n}\) and its dual group as in Corollary 3 of II, p. 236. Let N denote the Euclidean norm on \(\mathbf{R}^{n}\) and let \(\mathcal{F}\) be the Fourier transform on \(\mathcal{S}(\mathbf{R}^{n})\) and on \(\mathcal{S}'(\mathbf{R}^{n})\) (no. 12 of IV, p. 217).

Let \(\Delta_{\mathcal{S}}\) be the partial operator on \(\mathrm{L}^2 (\mathbf{R}^n)\) whose domain is the space \(\mathcal{S}(\mathbf{R}^n)\) and which satisfies

\[
\Delta_ {\mathcal {S}} (\varphi) = - \sum_ {i = 1} ^ {n} \partial_ {i} ^ {2} \varphi
\]

for all \(\varphi\in\mathcal{S}(\mathbf{R}^{n})\) . We then have \(\mathcal{F}(\Delta_{\mathcal{S}}(\varphi))=4\pi^{2}\mathrm{N}^{2}\mathcal{F}(\varphi)\) for every function \(\varphi\in\mathcal{S}(\mathbf{R}^{n})\) (remark 12 of IV, p. 220). Since the Fourier transform is an automorphism of \(\mathcal{S}(\mathbf{R}^{n})\) (prop. 18 of IV, p. 219), this means that the partial operator \(u=\mathcal{F}\circ\Delta_{\mathcal{S}}\circ\mathcal{F}^{-1}\) is the reduction of the multiplication operator \(m_{4\pi^{2}\mathrm{N}^{2}}\) on the space \(\mathcal{S}(\mathbf{R}^{n})\) . The partial operator \(u\) is closable and symmetric; its closure is the positive self-adjoint partial operator \(m_{4\pi^{2}\mathrm{N}^{2}}\) (prop. 6 of IV, p. 232, applied to the space \(\mathcal{D}(\mathbf{R}^{n})\subset\mathcal{S}(\mathbf{R}^{n})\) using prop. 4 of IV, p. 202). Consequently, the partial operator \(\Delta_{\mathcal{S}}\) is essentially self-adjoint. We denote \(\Delta\) its closure; it is a positive self-adjoint operator, and it is the unique Laplacian on \(\mathbf{R}^{n}\) ( \(cf.\) example from No. 6 of IV, p. 243 and corollary of Prop. 26 of IV, p. 257).

By prop. 21 of IV, p. 223, the Sobolev space \(\mathrm{H}^2 (\mathbf{R}^n)\) is the set of \(f\in \mathrm{L}^2 (\mathbf{R}^n)\) such that \((1 + \mathrm{N}^2)\mathcal{F}(f)\) belongs to \(\mathrm{L}^2 (\mathbf{R}^n)\), that is, such that \(\mathcal{F}(f)\) belongs to the domain of \(m_{4\pi^2\mathrm{N}^2}\) . Since the Fourier transform is an isometric isomorphism of the space \(\mathrm{L}^2 (\mathbf{R}^n)\) onto itself (th. 1 of II, p. 215), the domain of the Laplacian \(\Delta\) is \(\mathrm{H}^2 (\mathbf{R}^n)\) . We have \(\Delta = \mathcal{F}^{-1}\circ m_{4\pi^2\mathrm{N}^2}\circ \mathcal{F}\) . In particular, the spectrum of \(\Delta\) is equal to \(\mathbf{R}_{+}\) .